
\documentclass{article}
\usepackage{colortbl}
\usepackage{makecell}
\usepackage{multirow}
\usepackage{supertabular}

\begin{document}

\newcounter{utterance}

\centering \large Interaction Transcript for game `clean\_up', experiment `0\_easy\_5obj\_en', episode 1 with llama{-}3.3{-}70b{-}instruct{-}t1.0.
\vspace{24pt}

{ \footnotesize  \setcounter{utterance}{1}
\setlength{\tabcolsep}{0pt}
\begin{supertabular}{c@{$\;$}|p{.15\linewidth}@{}p{.15\linewidth}p{.15\linewidth}p{.15\linewidth}p{.15\linewidth}p{.15\linewidth}}
    \# & $\;$A & \multicolumn{4}{c}{Game Master} & $\;\:$B\\
    \hline

    \theutterance \stepcounter{utterance}  
    & & \multicolumn{4}{p{0.6\linewidth}}{
        \cellcolor[rgb]{0.9,0.9,0.9}{
            \makecell[{{p{\linewidth}}}]{
                \texttt{\tiny{[P1$\langle$GM]}}
                \texttt{Sono il tuo game master, e stai giocando un gioco collaborativo con la seguente griglia come tavolo da gioco:} \\
\\ 
\texttt{```} \\
\texttt{ 1234567} \\
\texttt{╔═══════╗} \\
\texttt{║I◌◌◌◌◌◌║ 1} \\
\texttt{║◌◌◌◌◌◌◌║ 2} \\
\texttt{║◌◌◌T◌┌─╢ 3} \\
\texttt{║◌◌◌┌─┤◌║ 4} \\
\texttt{║◌◌W│C├─╢ 5} \\
\texttt{╟───┼─┘H║ 6} \\
\texttt{║◌◌◌│◌◌◌║ 7} \\
\texttt{╚═══╧═══╝} \\
\\ 
\texttt{```} \\
\\ 
\texttt{* Il bordo superiore mostra le coordinate x che aumentano verso destra, e il bordo destro le coordinate y che aumentano verso il basso.} \\
\texttt{* The following objects are randomly placed on your grid: 'W', 'I', 'T', 'C', 'H'.} \\
\\ 
\\ 
\texttt{L'altro giocatore vede una variante del tavolo da gioco, dove gli oggetti sono posizionati in diverse posizioni casuali. Non puoi vedere il tavolo dell'altro giocatore, e loro non possono vedere il tuo.} \\
\\ 
\\ 
\texttt{**Obiettivo:**} \\
\\ 
\texttt{Entrambi i giocatori devono spostare gli oggetti sui rispettivi sfondi in modo che oggetti identici finiscano alle stesse coordinate. Devi comunicare con l'altro giocatore per concordare una configurazione obiettivo comune.} \\
\\ 
\\ 
\texttt{**Regole:**} \\
\\ 
\texttt{* In ogni turno, puoi inviare esattamente uno dei seguenti due comandi:} \\
\texttt{ 1. `COMUNICA: <MESSAGE>`: per inviare un messaggio (tutto fino alla prossima interruzione di riga) all'altro giocatore. Lo inoltrerò al tuo partner.} \\
\texttt{2. `SPOSTA: <OBJECT>, (<X>, <Y>)`: per spostare un oggetto in una nuova posizione, dove <X> è la colonna e <Y> è la riga. Ti farò sapere se la tua mossa è andata a buon fine o no.} \\
\texttt{* Se non rispetti il formato, o invii diversi comandi contemporaneamente, devo penalizzarti.} \\
\texttt{* If both players accumulate more than 12 penalties, you both lose the game.} \\
\texttt{* È vitale che comunichi con l'altro giocatore riguardo al tuo stato obiettivo! L'*unico* modo in cui puoi trasmettere la tua strategia all'altro giocatore è usando il comando `COMUNICA: <MESSAGE>`!} \\
\\ 
\\ 
\texttt{**Spostamento Oggetti**} \\
\\ 
\texttt{* Puoi spostare oggetti solo in celle entro i limiti della griglia. La cella di destinazione deve essere vuota, cioè deve contenere solo il simbolo '◌'.} \\
\texttt{* Se provi a spostare un oggetto in un posto che non è vuoto, o provi a spostarlo fuori dalla griglia, devo penalizzarti. Hai un altro tentativo.} \\
\texttt{* Prima di fare una mossa, controlla due volte che il posto di destinazione sia vuoto, e non contenga alcuna lettera, cornice, o linea!} \\
\\ 
\\ 
\texttt{**Fine del Gioco**} \\
\\ 
\texttt{Se pensi di aver raggiunto l'obiettivo di allineare tutti gli oggetti, puoi chiedere all'altro giocatore di finire il gioco inviando COMUNICA: finito?. Se l'altro giocatore ti chiede di finire il gioco, e rispondi con COMUNICA: finito!, the game will end.} \\
\\ 
\texttt{Both players win if the game ends within 20 rounds, where one round is defined as two players each sending a valid command.} \\
\\ 
\texttt{**Punteggio:**} \\
\\ 
\texttt{Più vicini sono gli oggetti identici in entrambi i tavoli da gioco, più punti ottieni. Le penalità riducono i tuoi punti. Riesci a battere il record?} \\
\\ 
\\ 
\texttt{Per favore invia un messaggio all'altro giocatore per iniziare il gioco!} \\
            }
        }
    }
    & \\ \\

    \theutterance \stepcounter{utterance}  
    & \multicolumn{4}{p{0.6\linewidth}}{
        \cellcolor[rgb]{0.9,0.9,0.9}{
            \makecell[{{p{\linewidth}}}]{
                \texttt{\tiny{[P1$\rangle$GM]}}
                \texttt{COMUNICA: Ciao! Sono pronto a iniziare il gioco. Propongo di iniziare descrivendo la posizione degli oggetti sul nostro tavolo da gioco. Io ho un oggetto 'I' alla posizione (1,1), un oggetto 'T' alla posizione (4,3), un oggetto 'W' alla posizione (1,5), un oggetto 'C' alla posizione (2,5) e un oggetto 'H' alla posizione (3,6). Come sono posizionati gli oggetti sul tuo tavolo da gioco?} \\
            }
        }
    }
    & & \\ \\

    \theutterance \stepcounter{utterance}  
    & & & \multicolumn{2}{p{0.3\linewidth}}{
        \cellcolor[rgb]{0.9,0.9,0.9}{
            \makecell[{{p{\linewidth}}}]{
                \texttt{\tiny{[GM$|$GM]}}
                \texttt{COMUNICA: Ciao! Sono pronto a iniziare il gioco. Propongo di iniziare descrivendo la posizione degli oggetti sul nostro tavolo da gioco. Io ho un oggetto 'I' alla posizione (1,1), un oggetto 'T' alla posizione (4,3), un oggetto 'W' alla posizione (1,5), un oggetto 'C' alla posizione (2,5) e un oggetto 'H' alla posizione (3,6). Come sono posizionati gli oggetti sul tuo tavolo da gioco?} \\
            }
        }
    }
    & & \\ \\

    \theutterance \stepcounter{utterance}  
    & & \multicolumn{4}{p{0.6\linewidth}}{
        \cellcolor[rgb]{0.9,0.9,0.9}{
            \makecell[{{p{\linewidth}}}]{
                \texttt{\tiny{[GM$\rangle$P2]}}
                \texttt{Sono il tuo game master, e stai giocando un gioco collaborativo con la seguente griglia come tavolo da gioco:} \\
\\ 
\texttt{```} \\
\texttt{ 1234567} \\
\texttt{╔═══════╗} \\
\texttt{║◌◌H◌◌◌◌║ 1} \\
\texttt{║◌◌◌C◌◌◌║ 2} \\
\texttt{║◌◌◌◌◌┌─╢ 3} \\
\texttt{║◌◌◌┌─┤◌║ 4} \\
\texttt{║◌IW│◌├─╢ 5} \\
\texttt{╟───┼─┘◌║ 6} \\
\texttt{║T◌◌│◌◌◌║ 7} \\
\texttt{╚═══╧═══╝} \\
\\ 
\texttt{```} \\
\\ 
\texttt{* Il bordo superiore mostra le coordinate x che aumentano verso destra, e il bordo destro le coordinate y che aumentano verso il basso.} \\
\texttt{* The following objects are randomly placed on your grid: 'W', 'I', 'T', 'C', 'H'.} \\
\\ 
\\ 
\texttt{L'altro giocatore vede una variante del tavolo da gioco, dove gli oggetti sono posizionati in diverse posizioni casuali. Non puoi vedere il tavolo dell'altro giocatore, e loro non possono vedere il tuo.} \\
\\ 
\\ 
\texttt{**Obiettivo:**} \\
\\ 
\texttt{Entrambi i giocatori devono spostare gli oggetti sui rispettivi sfondi in modo che oggetti identici finiscano alle stesse coordinate. Devi comunicare con l'altro giocatore per concordare una configurazione obiettivo comune.} \\
\\ 
\\ 
\texttt{**Regole:**} \\
\\ 
\texttt{* In ogni turno, puoi inviare esattamente uno dei seguenti due comandi:} \\
\texttt{ 1. `COMUNICA: <MESSAGE>`: per inviare un messaggio (tutto fino alla prossima interruzione di riga) all'altro giocatore. Lo inoltrerò al tuo partner.} \\
\texttt{2. `SPOSTA: <OBJECT>, (<X>, <Y>)`: per spostare un oggetto in una nuova posizione, dove <X> è la colonna e <Y> è la riga. Ti farò sapere se la tua mossa è andata a buon fine o no.} \\
\texttt{* Se non rispetti il formato, o invii diversi comandi contemporaneamente, devo penalizzarti.} \\
\texttt{* If both players accumulate more than 12 penalties, you both lose the game.} \\
\texttt{* È vitale che comunichi con l'altro giocatore riguardo al tuo stato obiettivo! L'*unico* modo in cui puoi trasmettere la tua strategia all'altro giocatore è usando il comando `COMUNICA: <MESSAGE>`!} \\
\\ 
\\ 
\texttt{**Spostamento Oggetti**} \\
\\ 
\texttt{* Puoi spostare oggetti solo in celle entro i limiti della griglia. La cella di destinazione deve essere vuota, cioè deve contenere solo il simbolo '◌'.} \\
\texttt{* Se provi a spostare un oggetto in un posto che non è vuoto, o provi a spostarlo fuori dalla griglia, devo penalizzarti. Hai un altro tentativo.} \\
\texttt{* Prima di fare una mossa, controlla due volte che il posto di destinazione sia vuoto, e non contenga alcuna lettera, cornice, o linea!} \\
\\ 
\\ 
\texttt{**Fine del Gioco**} \\
\\ 
\texttt{Se pensi di aver raggiunto l'obiettivo di allineare tutti gli oggetti, puoi chiedere all'altro giocatore di finire il gioco inviando COMUNICA: finito?. Se l'altro giocatore ti chiede di finire il gioco, e rispondi con COMUNICA: finito!, the game will end.} \\
\\ 
\texttt{Both players win if the game ends within 20 rounds, where one round is defined as two players each sending a valid command.} \\
\\ 
\texttt{**Punteggio:**} \\
\\ 
\texttt{Più vicini sono gli oggetti identici in entrambi i tavoli da gioco, più punti ottieni. Le penalità riducono i tuoi punti. Riesci a battere il record?} \\
\\ 
\texttt{L'altro giocatore ha iniziato il gioco inviando questo messaggio:} \\
\\ 
\texttt{"Ciao! Sono pronto a iniziare il gioco. Propongo di iniziare descrivendo la posizione degli oggetti sul nostro tavolo da gioco. Io ho un oggetto 'I' alla posizione (1,1), un oggetto 'T' alla posizione (4,3), un oggetto 'W' alla posizione (1,5), un oggetto 'C' alla posizione (2,5) e un oggetto 'H' alla posizione (3,6). Come sono posizionati gli oggetti sul tuo tavolo da gioco?"} \\
\\ 
\texttt{Qual è il tuo primo comando?} \\
            }
        }
    }
    & \\ \\

    \theutterance \stepcounter{utterance}  
    & & & \multicolumn{4}{p{0.6\linewidth}}{
        \cellcolor[rgb]{0.9,0.9,0.9}{
            \makecell[{{p{\linewidth}}}]{
                \texttt{\tiny{[GM$\langle$P2]}}
                \texttt{COMUNICA: Ciao! Sono pronto a iniziare il gioco. La posizione degli oggetti sul mio tavolo da gioco è la seguente: 'H' alla posizione (3,1), 'C' alla posizione (4,2), 'I' alla posizione (1,5), 'W' alla posizione (2,5) e 'T' alla posizione (1,7). Propongo di procedere descrivendo le mosse che possiamo fare per allineare gli oggetti identici. Qual è la tua prima mossa?} \\
            }
        }
    }
     \\ \\

    \theutterance \stepcounter{utterance}  
    & & & \multicolumn{2}{p{0.3\linewidth}}{
        \cellcolor[rgb]{0.9,0.9,0.9}{
            \makecell[{{p{\linewidth}}}]{
                \texttt{\tiny{[GM$|$GM]}}
                \texttt{COMUNICA: Ciao! Sono pronto a iniziare il gioco. La posizione degli oggetti sul mio tavolo da gioco è la seguente: 'H' alla posizione (3,1), 'C' alla posizione (4,2), 'I' alla posizione (1,5), 'W' alla posizione (2,5) e 'T' alla posizione (1,7). Propongo di procedere descrivendo le mosse che possiamo fare per allineare gli oggetti identici. Qual è la tua prima mossa?} \\
            }
        }
    }
    & & \\ \\

    \theutterance \stepcounter{utterance}  
    & & \multicolumn{4}{p{0.6\linewidth}}{
        \cellcolor[rgb]{0.9,0.9,0.9}{
            \makecell[{{p{\linewidth}}}]{
                \texttt{\tiny{[P1$\langle$GM]}}
                \texttt{Il tuo messaggio è stato inoltrato all'altro giocatore.} \\
\\ 
\\ 
\texttt{You are currently playing round 1 of maximum 20.} \\
\texttt{You have collectively accumulated 0 of 12 penalties.} \\
\texttt{L'altro giocatore ha inviato questo messaggio:} \\
\\ 
\\ 
\texttt{"Ciao! Sono pronto a iniziare il gioco. La posizione degli oggetti sul mio tavolo da gioco è la seguente: 'H' alla posizione (3,1), 'C' alla posizione (4,2), 'I' alla posizione (1,5), 'W' alla posizione (2,5) e 'T' alla posizione (1,7). Propongo di procedere descrivendo le mosse che possiamo fare per allineare gli oggetti identici. Qual è la tua prima mossa?"} \\
\\ 
\\ 
\texttt{Qual è il tuo prossimo comando?} \\
            }
        }
    }
    & \\ \\

    \theutterance \stepcounter{utterance}  
    & \multicolumn{4}{p{0.6\linewidth}}{
        \cellcolor[rgb]{0.9,0.9,0.9}{
            \makecell[{{p{\linewidth}}}]{
                \texttt{\tiny{[P1$\rangle$GM]}}
                \texttt{COMUNICA: Grazie per la descrizione della tua griglia! Notiamo che i nostri oggetti 'I' e 'W' sono già nella stessa riga (5) anche se non nella stessa colonna. Per iniziare, potremmo cercare di allineare gli oggetti 'H' e 'T' che sembrano essere i più lontani dalle rispettive posizioni comuni. Io proporrei di muovere il mio 'H' verso la parte superiore del tavolo, idealmente in una posizione che potrebbe essere vicina alla tua 'H'. Come pensi di muovere il tuo 'T' e il tuo 'H' per avvicinarli alle posizioni che potremmo decidere come comuni?} \\
            }
        }
    }
    & & \\ \\

    \theutterance \stepcounter{utterance}  
    & & & \multicolumn{2}{p{0.3\linewidth}}{
        \cellcolor[rgb]{0.9,0.9,0.9}{
            \makecell[{{p{\linewidth}}}]{
                \texttt{\tiny{[GM$|$GM]}}
                \texttt{COMUNICA: Grazie per la descrizione della tua griglia! Notiamo che i nostri oggetti 'I' e 'W' sono già nella stessa riga (5) anche se non nella stessa colonna. Per iniziare, potremmo cercare di allineare gli oggetti 'H' e 'T' che sembrano essere i più lontani dalle rispettive posizioni comuni. Io proporrei di muovere il mio 'H' verso la parte superiore del tavolo, idealmente in una posizione che potrebbe essere vicina alla tua 'H'. Come pensi di muovere il tuo 'T' e il tuo 'H' per avvicinarli alle posizioni che potremmo decidere come comuni?} \\
            }
        }
    }
    & & \\ \\

    \theutterance \stepcounter{utterance}  
    & & \multicolumn{4}{p{0.6\linewidth}}{
        \cellcolor[rgb]{0.9,0.9,0.9}{
            \makecell[{{p{\linewidth}}}]{
                \texttt{\tiny{[GM$\rangle$P2]}}
                \texttt{Il tuo messaggio è stato inoltrato all'altro giocatore.} \\
\\ 
\\ 
\texttt{You are currently playing round 2 of maximum 20.} \\
\texttt{You have collectively accumulated 0 of 12 penalties.} \\
\texttt{L'altro giocatore ha inviato questo messaggio:} \\
\\ 
\\ 
\texttt{"Grazie per la descrizione della tua griglia! Notiamo che i nostri oggetti 'I' e 'W' sono già nella stessa riga (5) anche se non nella stessa colonna. Per iniziare, potremmo cercare di allineare gli oggetti 'H' e 'T' che sembrano essere i più lontani dalle rispettive posizioni comuni. Io proporrei di muovere il mio 'H' verso la parte superiore del tavolo, idealmente in una posizione che potrebbe essere vicina alla tua 'H'. Come pensi di muovere il tuo 'T' e il tuo 'H' per avvicinarli alle posizioni che potremmo decidere come comuni?"} \\
\\ 
\\ 
\texttt{Qual è il tuo prossimo comando?} \\
            }
        }
    }
    & \\ \\

    \theutterance \stepcounter{utterance}  
    & & & \multicolumn{4}{p{0.6\linewidth}}{
        \cellcolor[rgb]{0.9,0.9,0.9}{
            \makecell[{{p{\linewidth}}}]{
                \texttt{\tiny{[GM$\langle$P2]}}
                \texttt{COMUNICA: Sono d'accordo che allineare gli oggetti 'H' e 'T' potrebbe essere un buon punto di partenza. Considerando la mia griglia, potrei muovere il mio 'T' verso la parte superiore del tavolo, forse nella colonna 3, in modo da non bloccare la colonna 4 dove si trova il mio 'C'. Per quanto riguarda il mio 'H', potrei lasciarlo nella sua posizione attuale (3,1) e vedere come si sviluppa la situazione. Sarebbe utile sapere come pensi di muovere il tuo 'H' e se hai in mente una posizione specifica per gli oggetti 'H' e 'T' che potremmo considerare come obiettivo comune. Inoltre, possiamo anche discutere la posizione degli oggetti 'I' e 'W' e come possiamo allinearli meglio. Qual è la tua idea su come procedere con gli oggetti 'I' e 'W'?} \\
            }
        }
    }
     \\ \\

    \theutterance \stepcounter{utterance}  
    & & & \multicolumn{2}{p{0.3\linewidth}}{
        \cellcolor[rgb]{0.9,0.9,0.9}{
            \makecell[{{p{\linewidth}}}]{
                \texttt{\tiny{[GM$|$GM]}}
                \texttt{COMUNICA: Sono d'accordo che allineare gli oggetti 'H' e 'T' potrebbe essere un buon punto di partenza. Considerando la mia griglia, potrei muovere il mio 'T' verso la parte superiore del tavolo, forse nella colonna 3, in modo da non bloccare la colonna 4 dove si trova il mio 'C'. Per quanto riguarda il mio 'H', potrei lasciarlo nella sua posizione attuale (3,1) e vedere come si sviluppa la situazione. Sarebbe utile sapere come pensi di muovere il tuo 'H' e se hai in mente una posizione specifica per gli oggetti 'H' e 'T' che potremmo considerare come obiettivo comune. Inoltre, possiamo anche discutere la posizione degli oggetti 'I' e 'W' e come possiamo allinearli meglio. Qual è la tua idea su come procedere con gli oggetti 'I' e 'W'?} \\
            }
        }
    }
    & & \\ \\

    \theutterance \stepcounter{utterance}  
    & & \multicolumn{4}{p{0.6\linewidth}}{
        \cellcolor[rgb]{0.9,0.9,0.9}{
            \makecell[{{p{\linewidth}}}]{
                \texttt{\tiny{[P1$\langle$GM]}}
                \texttt{Il tuo messaggio è stato inoltrato all'altro giocatore.} \\
\\ 
\\ 
\texttt{You are currently playing round 2 of maximum 20.} \\
\texttt{You have collectively accumulated 0 of 12 penalties.} \\
\texttt{L'altro giocatore ha inviato questo messaggio:} \\
\\ 
\\ 
\texttt{"Sono d'accordo che allineare gli oggetti 'H' e 'T' potrebbe essere un buon punto di partenza. Considerando la mia griglia, potrei muovere il mio 'T' verso la parte superiore del tavolo, forse nella colonna 3, in modo da non bloccare la colonna 4 dove si trova il mio 'C'. Per quanto riguarda il mio 'H', potrei lasciarlo nella sua posizione attuale (3,1) e vedere come si sviluppa la situazione. Sarebbe utile sapere come pensi di muovere il tuo 'H' e se hai in mente una posizione specifica per gli oggetti 'H' e 'T' che potremmo considerare come obiettivo comune. Inoltre, possiamo anche discutere la posizione degli oggetti 'I' e 'W' e come possiamo allinearli meglio. Qual è la tua idea su come procedere con gli oggetti 'I' e 'W'?"} \\
\\ 
\\ 
\texttt{Qual è il tuo prossimo comando?} \\
            }
        }
    }
    & \\ \\

    \theutterance \stepcounter{utterance}  
    & \multicolumn{4}{p{0.6\linewidth}}{
        \cellcolor[rgb]{0.9,0.9,0.9}{
            \makecell[{{p{\linewidth}}}]{
                \texttt{\tiny{[P1$\rangle$GM]}}
                \texttt{COMUNICA: Grazie per la risposta! Sono d'accordo che muovere il tuo 'T' nella colonna 3 potrebbe essere una mossa intelligente. Per quanto riguarda il mio 'H', penso di poterlo spostare nella posizione (3,1) per allineararlo con il tuo 'H'. Questo potrebbe essere un buon punto di partenza per gli oggetti 'H'. } \\
\\ 
\texttt{Per gli oggetti 'I' e 'W', noto che i nostri 'I' e 'W' sono già nella stessa riga (5) ma in colonne diverse. Potremmo cercare di spostarli in una posizione più centrale, magari intorno alla colonna 3 o 4, in modo da avere più flessibilità per le mosse future. } \\
\\ 
\texttt{Prima di procedere con le mosse, vorrei confermare: ti sembra ragionevole provare a posizionare gli oggetti 'H' intorno alla posizione (3,1), e poi discutere ulteriormente la posizione degli oggetti 'T', 'I' e 'W' in base ai progressi che facciamo?} \\
            }
        }
    }
    & & \\ \\

    \theutterance \stepcounter{utterance}  
    & & & \multicolumn{2}{p{0.3\linewidth}}{
        \cellcolor[rgb]{0.9,0.9,0.9}{
            \makecell[{{p{\linewidth}}}]{
                \texttt{\tiny{[GM$|$GM]}}
                \texttt{COMUNICA: Grazie per la risposta! Sono d'accordo che muovere il tuo 'T' nella colonna 3 potrebbe essere una mossa intelligente. Per quanto riguarda il mio 'H', penso di poterlo spostare nella posizione (3,1) per allineararlo con il tuo 'H'. Questo potrebbe essere un buon punto di partenza per gli oggetti 'H'. } \\
\\ 
\texttt{Per gli oggetti 'I' e 'W', noto che i nostri 'I' e 'W' sono già nella stessa riga (5) ma in colonne diverse. Potremmo cercare di spostarli in una posizione più centrale, magari intorno alla colonna 3 o 4, in modo da avere più flessibilità per le mosse future. } \\
\\ 
\texttt{Prima di procedere con le mosse, vorrei confermare: ti sembra ragionevole provare a posizionare gli oggetti 'H' intorno alla posizione (3,1), e poi discutere ulteriormente la posizione degli oggetti 'T', 'I' e 'W' in base ai progressi che facciamo?} \\
            }
        }
    }
    & & \\ \\

    \theutterance \stepcounter{utterance}  
    & & & \multicolumn{2}{p{0.3\linewidth}}{
        \cellcolor[rgb]{0.9,0.9,0.9}{
            \makecell[{{p{\linewidth}}}]{
                \texttt{\tiny{[GM$|$GM]}}
                \texttt{Invalid format: tail is not empty: '} \\
\\ 
\texttt{Per gli oggetti 'I' e 'W', noto che i nostri 'I' e 'W' sono già nella stessa riga (5) ma in colonne diverse. Potremmo cercare di spostarli in una posizione più centrale, magari intorno alla colonna 3 o 4, in modo da avere più flessibilità per le mosse future. } \\
\\ 
\texttt{Prima di procedere con le mosse, vorrei confermare: ti sembra ragionevole provare a posizionare gli oggetti 'H' intorno alla posizione (3,1), e poi discutere ulteriormente la posizione degli oggetti 'T', 'I' e 'W' in base ai progressi che facciamo?} \\
            }
        }
    }
    & & \\ \\

    \theutterance \stepcounter{utterance}  
    & & \multicolumn{4}{p{0.6\linewidth}}{
        \cellcolor[rgb]{0.9,0.9,0.9}{
            \makecell[{{p{\linewidth}}}]{
                \texttt{\tiny{[P1$\langle$GM]}}
                \texttt{Penalità: Il tuo messaggio non deve contenere nulla dopo il comando!} \\
\\ 
\texttt{Assicurati che la tua risposta contenga solo `COMUNICA: <MESSAGE>` o `SPOSTA: <OBJECT>, (<X>, <Y>)`, e nient'altro!} \\
\\ 
\texttt{You have collectively accumulated 1 of 12 penalties.} \\
\texttt{Riprova!} \\
            }
        }
    }
    & \\ \\

    \theutterance \stepcounter{utterance}  
    & \multicolumn{4}{p{0.6\linewidth}}{
        \cellcolor[rgb]{0.9,0.9,0.9}{
            \makecell[{{p{\linewidth}}}]{
                \texttt{\tiny{[P1$\rangle$GM]}}
                \texttt{COMUNICA: Grazie per la risposta! Sono d'accordo che muovere il tuo 'T' nella colonna 3 potrebbe essere una mossa intelligente. Per quanto riguarda il mio 'H', penso di poterlo spostare nella posizione (3,1) per allineararlo con il tuo 'H'. Questo potrebbe essere un buon punto di partenza per gli oggetti 'H'. Per gli oggetti 'I' e 'W', noto che i nostri 'I' e 'W' sono già nella stessa riga (5) ma in colonne diverse. Potremmo cercare di spostarli in una posizione più centrale, magari intorno alla colonna 3 o 4, in modo da avere più flessibilità per le mosse future. Ti sembra ragionevole provare a posizionare gli oggetti 'H' intorno alla posizione (3,1), e poi discutere ulteriormente la posizione degli oggetti 'T', 'I' e 'W' in base ai progressi che facciamo?} \\
            }
        }
    }
    & & \\ \\

    \theutterance \stepcounter{utterance}  
    & & & \multicolumn{2}{p{0.3\linewidth}}{
        \cellcolor[rgb]{0.9,0.9,0.9}{
            \makecell[{{p{\linewidth}}}]{
                \texttt{\tiny{[GM$|$GM]}}
                \texttt{COMUNICA: Grazie per la risposta! Sono d'accordo che muovere il tuo 'T' nella colonna 3 potrebbe essere una mossa intelligente. Per quanto riguarda il mio 'H', penso di poterlo spostare nella posizione (3,1) per allineararlo con il tuo 'H'. Questo potrebbe essere un buon punto di partenza per gli oggetti 'H'. Per gli oggetti 'I' e 'W', noto che i nostri 'I' e 'W' sono già nella stessa riga (5) ma in colonne diverse. Potremmo cercare di spostarli in una posizione più centrale, magari intorno alla colonna 3 o 4, in modo da avere più flessibilità per le mosse future. Ti sembra ragionevole provare a posizionare gli oggetti 'H' intorno alla posizione (3,1), e poi discutere ulteriormente la posizione degli oggetti 'T', 'I' e 'W' in base ai progressi che facciamo?} \\
            }
        }
    }
    & & \\ \\

    \theutterance \stepcounter{utterance}  
    & & \multicolumn{4}{p{0.6\linewidth}}{
        \cellcolor[rgb]{0.9,0.9,0.9}{
            \makecell[{{p{\linewidth}}}]{
                \texttt{\tiny{[GM$\rangle$P2]}}
                \texttt{Il tuo messaggio è stato inoltrato all'altro giocatore.} \\
\\ 
\\ 
\texttt{You are currently playing round 3 of maximum 20.} \\
\texttt{You have collectively accumulated 1 of 12 penalties.} \\
\texttt{L'altro giocatore ha inviato questo messaggio:} \\
\\ 
\\ 
\texttt{"Grazie per la risposta! Sono d'accordo che muovere il tuo 'T' nella colonna 3 potrebbe essere una mossa intelligente. Per quanto riguarda il mio 'H', penso di poterlo spostare nella posizione (3,1) per allineararlo con il tuo 'H'. Questo potrebbe essere un buon punto di partenza per gli oggetti 'H'. Per gli oggetti 'I' e 'W', noto che i nostri 'I' e 'W' sono già nella stessa riga (5) ma in colonne diverse. Potremmo cercare di spostarli in una posizione più centrale, magari intorno alla colonna 3 o 4, in modo da avere più flessibilità per le mosse future. Ti sembra ragionevole provare a posizionare gli oggetti 'H' intorno alla posizione (3,1), e poi discutere ulteriormente la posizione degli oggetti 'T', 'I' e 'W' in base ai progressi che facciamo?"} \\
\\ 
\\ 
\texttt{Qual è il tuo prossimo comando?} \\
            }
        }
    }
    & \\ \\

    \theutterance \stepcounter{utterance}  
    & & & \multicolumn{4}{p{0.6\linewidth}}{
        \cellcolor[rgb]{0.9,0.9,0.9}{
            \makecell[{{p{\linewidth}}}]{
                \texttt{\tiny{[GM$\langle$P2]}}
                \texttt{COMUNICA: Sono d'accordo che spostare il tuo 'H' nella posizione (3,1) potrebbe essere un buon punto di partenza per allineare gli oggetti 'H'. La mia idea è che potremmo considerare la posizione (3,1) come un punto di riferimento per gli oggetti 'H'. Per gli oggetti 'I' e 'W', sono d'accordo che spostarli in una posizione più centrale potrebbe essere utile. Tuttavia, prima di decidere esattamente dove spostarli, vorrei chiarire la posizione degli oggetti 'T'. Dal momento che ho intenzione di spostare il mio 'T' nella colonna 3, potresti considerare di spostare il tuo 'T' nella stessa colonna, in modo da allinearli? Inoltre, potremmo discutere la possibilità di stabilire un ordine per le mosse, ad esempio, potremmo iniziare a spostare gli oggetti 'H' e poi procedere con gli oggetti 'T', 'I' e 'W'. Come pensi di procedere con gli oggetti 'T' e con l'ordine delle mosse?} \\
            }
        }
    }
     \\ \\

    \theutterance \stepcounter{utterance}  
    & & & \multicolumn{2}{p{0.3\linewidth}}{
        \cellcolor[rgb]{0.9,0.9,0.9}{
            \makecell[{{p{\linewidth}}}]{
                \texttt{\tiny{[GM$|$GM]}}
                \texttt{COMUNICA: Sono d'accordo che spostare il tuo 'H' nella posizione (3,1) potrebbe essere un buon punto di partenza per allineare gli oggetti 'H'. La mia idea è che potremmo considerare la posizione (3,1) come un punto di riferimento per gli oggetti 'H'. Per gli oggetti 'I' e 'W', sono d'accordo che spostarli in una posizione più centrale potrebbe essere utile. Tuttavia, prima di decidere esattamente dove spostarli, vorrei chiarire la posizione degli oggetti 'T'. Dal momento che ho intenzione di spostare il mio 'T' nella colonna 3, potresti considerare di spostare il tuo 'T' nella stessa colonna, in modo da allinearli? Inoltre, potremmo discutere la possibilità di stabilire un ordine per le mosse, ad esempio, potremmo iniziare a spostare gli oggetti 'H' e poi procedere con gli oggetti 'T', 'I' e 'W'. Come pensi di procedere con gli oggetti 'T' e con l'ordine delle mosse?} \\
            }
        }
    }
    & & \\ \\

    \theutterance \stepcounter{utterance}  
    & & \multicolumn{4}{p{0.6\linewidth}}{
        \cellcolor[rgb]{0.9,0.9,0.9}{
            \makecell[{{p{\linewidth}}}]{
                \texttt{\tiny{[P1$\langle$GM]}}
                \texttt{Il tuo messaggio è stato inoltrato all'altro giocatore.} \\
\\ 
\\ 
\texttt{You are currently playing round 3 of maximum 20.} \\
\texttt{You have collectively accumulated 1 of 12 penalties.} \\
\texttt{L'altro giocatore ha inviato questo messaggio:} \\
\\ 
\\ 
\texttt{"Sono d'accordo che spostare il tuo 'H' nella posizione (3,1) potrebbe essere un buon punto di partenza per allineare gli oggetti 'H'. La mia idea è che potremmo considerare la posizione (3,1) come un punto di riferimento per gli oggetti 'H'. Per gli oggetti 'I' e 'W', sono d'accordo che spostarli in una posizione più centrale potrebbe essere utile. Tuttavia, prima di decidere esattamente dove spostarli, vorrei chiarire la posizione degli oggetti 'T'. Dal momento che ho intenzione di spostare il mio 'T' nella colonna 3, potresti considerare di spostare il tuo 'T' nella stessa colonna, in modo da allinearli? Inoltre, potremmo discutere la possibilità di stabilire un ordine per le mosse, ad esempio, potremmo iniziare a spostare gli oggetti 'H' e poi procedere con gli oggetti 'T', 'I' e 'W'. Come pensi di procedere con gli oggetti 'T' e con l'ordine delle mosse?"} \\
\\ 
\\ 
\texttt{Qual è il tuo prossimo comando?} \\
            }
        }
    }
    & \\ \\

    \theutterance \stepcounter{utterance}  
    & \multicolumn{4}{p{0.6\linewidth}}{
        \cellcolor[rgb]{0.9,0.9,0.9}{
            \makecell[{{p{\linewidth}}}]{
                \texttt{\tiny{[P1$\rangle$GM]}}
                \texttt{COMUNICA: Sono d'accordo che allineare gli oggetti 'T' nella colonna 3 potrebbe essere una buona idea. Considererò di spostare il mio 'T' nella colonna 3 per allinearli con il tuo 'T'. Per quanto riguarda l'ordine delle mosse, penso che iniziare con gli oggetti 'H' e poi procedere con gli oggetti 'T', 'I' e 'W' sia un piano ragionevole. Tuttavia, prima di iniziare a spostare gli oggetti, vorrei confermare se siamo d'accordo sulla posizione finale per gli oggetti 'H' e 'T'. Potremmo considerare di stabilire una posizione esatta per gli oggetti 'H' e 'T' e poi discutere le mosse necessarie per raggiungere quella posizione. Pensieri?} \\
            }
        }
    }
    & & \\ \\

    \theutterance \stepcounter{utterance}  
    & & & \multicolumn{2}{p{0.3\linewidth}}{
        \cellcolor[rgb]{0.9,0.9,0.9}{
            \makecell[{{p{\linewidth}}}]{
                \texttt{\tiny{[GM$|$GM]}}
                \texttt{COMUNICA: Sono d'accordo che allineare gli oggetti 'T' nella colonna 3 potrebbe essere una buona idea. Considererò di spostare il mio 'T' nella colonna 3 per allinearli con il tuo 'T'. Per quanto riguarda l'ordine delle mosse, penso che iniziare con gli oggetti 'H' e poi procedere con gli oggetti 'T', 'I' e 'W' sia un piano ragionevole. Tuttavia, prima di iniziare a spostare gli oggetti, vorrei confermare se siamo d'accordo sulla posizione finale per gli oggetti 'H' e 'T'. Potremmo considerare di stabilire una posizione esatta per gli oggetti 'H' e 'T' e poi discutere le mosse necessarie per raggiungere quella posizione. Pensieri?} \\
            }
        }
    }
    & & \\ \\

    \theutterance \stepcounter{utterance}  
    & & \multicolumn{4}{p{0.6\linewidth}}{
        \cellcolor[rgb]{0.9,0.9,0.9}{
            \makecell[{{p{\linewidth}}}]{
                \texttt{\tiny{[GM$\rangle$P2]}}
                \texttt{Il tuo messaggio è stato inoltrato all'altro giocatore.} \\
\\ 
\\ 
\texttt{You are currently playing round 4 of maximum 20.} \\
\texttt{You have collectively accumulated 1 of 12 penalties.} \\
\texttt{L'altro giocatore ha inviato questo messaggio:} \\
\\ 
\\ 
\texttt{"Sono d'accordo che allineare gli oggetti 'T' nella colonna 3 potrebbe essere una buona idea. Considererò di spostare il mio 'T' nella colonna 3 per allinearli con il tuo 'T'. Per quanto riguarda l'ordine delle mosse, penso che iniziare con gli oggetti 'H' e poi procedere con gli oggetti 'T', 'I' e 'W' sia un piano ragionevole. Tuttavia, prima di iniziare a spostare gli oggetti, vorrei confermare se siamo d'accordo sulla posizione finale per gli oggetti 'H' e 'T'. Potremmo considerare di stabilire una posizione esatta per gli oggetti 'H' e 'T' e poi discutere le mosse necessarie per raggiungere quella posizione. Pensieri?"} \\
\\ 
\\ 
\texttt{Qual è il tuo prossimo comando?} \\
            }
        }
    }
    & \\ \\

    \theutterance \stepcounter{utterance}  
    & & & \multicolumn{4}{p{0.6\linewidth}}{
        \cellcolor[rgb]{0.9,0.9,0.9}{
            \makecell[{{p{\linewidth}}}]{
                \texttt{\tiny{[GM$\langle$P2]}}
                \texttt{COMUNICA: Sono d'accordo che stabilire una posizione finale per gli oggetti 'H' e 'T' sia un passo importante prima di iniziare a spostarli. Considerando la nostra discussione precedente, potremmo considerare di posizionare gli oggetti 'H' nella posizione (3,1) e gli oggetti 'T' nella posizione (3,7). Questo ci darebbe un punto di riferimento chiaro per gli oggetti 'H' e 'T' e ci permetterebbe di pianificare le mosse necessarie per raggiungere quella posizione. Inoltre, potremmo anche discutere la possibilità di stabilire una posizione finale per gli oggetti 'I' e 'W' in modo da avere un piano completo per tutti gli oggetti. Come pensi di procedere con la posizione degli oggetti 'I' e 'W' e con le mosse necessarie per raggiungere le posizioni finali per tutti gli oggetti? } \\
\\ 
\texttt{Dopo aver inviato questo messaggio, potrei iniziare a pianificare la mia prima mossa per spostare il mio 'T' nella posizione (3,7), cercherò di eseguire la mossa nel prossimo turno.} \\
            }
        }
    }
     \\ \\

    \theutterance \stepcounter{utterance}  
    & & & \multicolumn{2}{p{0.3\linewidth}}{
        \cellcolor[rgb]{0.9,0.9,0.9}{
            \makecell[{{p{\linewidth}}}]{
                \texttt{\tiny{[GM$|$GM]}}
                \texttt{COMUNICA: Sono d'accordo che stabilire una posizione finale per gli oggetti 'H' e 'T' sia un passo importante prima di iniziare a spostarli. Considerando la nostra discussione precedente, potremmo considerare di posizionare gli oggetti 'H' nella posizione (3,1) e gli oggetti 'T' nella posizione (3,7). Questo ci darebbe un punto di riferimento chiaro per gli oggetti 'H' e 'T' e ci permetterebbe di pianificare le mosse necessarie per raggiungere quella posizione. Inoltre, potremmo anche discutere la possibilità di stabilire una posizione finale per gli oggetti 'I' e 'W' in modo da avere un piano completo per tutti gli oggetti. Come pensi di procedere con la posizione degli oggetti 'I' e 'W' e con le mosse necessarie per raggiungere le posizioni finali per tutti gli oggetti? } \\
\\ 
\texttt{Dopo aver inviato questo messaggio, potrei iniziare a pianificare la mia prima mossa per spostare il mio 'T' nella posizione (3,7), cercherò di eseguire la mossa nel prossimo turno.} \\
            }
        }
    }
    & & \\ \\

    \theutterance \stepcounter{utterance}  
    & & & \multicolumn{2}{p{0.3\linewidth}}{
        \cellcolor[rgb]{0.9,0.9,0.9}{
            \makecell[{{p{\linewidth}}}]{
                \texttt{\tiny{[GM$|$GM]}}
                \texttt{Invalid format: tail is not empty: '} \\
\\ 
\texttt{Dopo aver inviato questo messaggio, potrei iniziare a pianificare la mia prima mossa per spostare il mio 'T' nella posizione (3,7), cercherò di eseguire la mossa nel prossimo turno.} \\
            }
        }
    }
    & & \\ \\

    \theutterance \stepcounter{utterance}  
    & & \multicolumn{4}{p{0.6\linewidth}}{
        \cellcolor[rgb]{0.9,0.9,0.9}{
            \makecell[{{p{\linewidth}}}]{
                \texttt{\tiny{[GM$\rangle$P2]}}
                \texttt{Penalità: Il tuo messaggio non deve contenere nulla dopo il comando!} \\
\\ 
\texttt{Assicurati che la tua risposta contenga solo `COMUNICA: <MESSAGE>` o `SPOSTA: <OBJECT>, (<X>, <Y>)`, e nient'altro!} \\
\\ 
\texttt{You have collectively accumulated 2 of 12 penalties.} \\
\texttt{Riprova!} \\
            }
        }
    }
    & \\ \\

    \theutterance \stepcounter{utterance}  
    & & & \multicolumn{4}{p{0.6\linewidth}}{
        \cellcolor[rgb]{0.9,0.9,0.9}{
            \makecell[{{p{\linewidth}}}]{
                \texttt{\tiny{[GM$\langle$P2]}}
                \texttt{COMUNICA: Sono d'accordo che stabilire una posizione finale per gli oggetti 'H' e 'T' sia un passo importante prima di iniziare a spostarli. Considerando la nostra discussione precedente, potremmo considerare di posizionare gli oggetti 'H' nella posizione (3,1) e gli oggetti 'T' nella posizione (3,7). Questo ci darebbe un punto di riferimento chiaro per gli oggetti 'H' e 'T' e ci permetterebbe di pianificare le mosse necessarie per raggiungere quella posizione. Inoltre, potremmo anche discutere la possibilità di stabilire una posizione finale per gli oggetti 'I' e 'W' in modo da avere un piano completo per tutti gli oggetti. Come pensi di procedere con la posizione degli oggetti 'I' e 'W' e con le mosse necessarie per raggiungere le posizioni finali per tutti gli oggetti?} \\
            }
        }
    }
     \\ \\

    \theutterance \stepcounter{utterance}  
    & & & \multicolumn{2}{p{0.3\linewidth}}{
        \cellcolor[rgb]{0.9,0.9,0.9}{
            \makecell[{{p{\linewidth}}}]{
                \texttt{\tiny{[GM$|$GM]}}
                \texttt{COMUNICA: Sono d'accordo che stabilire una posizione finale per gli oggetti 'H' e 'T' sia un passo importante prima di iniziare a spostarli. Considerando la nostra discussione precedente, potremmo considerare di posizionare gli oggetti 'H' nella posizione (3,1) e gli oggetti 'T' nella posizione (3,7). Questo ci darebbe un punto di riferimento chiaro per gli oggetti 'H' e 'T' e ci permetterebbe di pianificare le mosse necessarie per raggiungere quella posizione. Inoltre, potremmo anche discutere la possibilità di stabilire una posizione finale per gli oggetti 'I' e 'W' in modo da avere un piano completo per tutti gli oggetti. Come pensi di procedere con la posizione degli oggetti 'I' e 'W' e con le mosse necessarie per raggiungere le posizioni finali per tutti gli oggetti?} \\
            }
        }
    }
    & & \\ \\

    \theutterance \stepcounter{utterance}  
    & & \multicolumn{4}{p{0.6\linewidth}}{
        \cellcolor[rgb]{0.9,0.9,0.9}{
            \makecell[{{p{\linewidth}}}]{
                \texttt{\tiny{[P1$\langle$GM]}}
                \texttt{Il tuo messaggio è stato inoltrato all'altro giocatore.} \\
\\ 
\\ 
\texttt{You are currently playing round 4 of maximum 20.} \\
\texttt{You have collectively accumulated 2 of 12 penalties.} \\
\texttt{L'altro giocatore ha inviato questo messaggio:} \\
\\ 
\\ 
\texttt{"Sono d'accordo che stabilire una posizione finale per gli oggetti 'H' e 'T' sia un passo importante prima di iniziare a spostarli. Considerando la nostra discussione precedente, potremmo considerare di posizionare gli oggetti 'H' nella posizione (3,1) e gli oggetti 'T' nella posizione (3,7). Questo ci darebbe un punto di riferimento chiaro per gli oggetti 'H' e 'T' e ci permetterebbe di pianificare le mosse necessarie per raggiungere quella posizione. Inoltre, potremmo anche discutere la possibilità di stabilire una posizione finale per gli oggetti 'I' e 'W' in modo da avere un piano completo per tutti gli oggetti. Come pensi di procedere con la posizione degli oggetti 'I' e 'W' e con le mosse necessarie per raggiungere le posizioni finali per tutti gli oggetti?"} \\
\\ 
\\ 
\texttt{Qual è il tuo prossimo comando?} \\
            }
        }
    }
    & \\ \\

    \theutterance \stepcounter{utterance}  
    & \multicolumn{4}{p{0.6\linewidth}}{
        \cellcolor[rgb]{0.9,0.9,0.9}{
            \makecell[{{p{\linewidth}}}]{
                \texttt{\tiny{[P1$\rangle$GM]}}
                \texttt{COMUNICA: Sono d'accordo con la posizione finale per gli oggetti 'H' e 'T' che hai proposto, (3,1) per 'H' e (3,7) per 'T'. Per gli oggetti 'I' e 'W', potremmo considerare di posizionarli nelle colonne 2 e 4, in modo da avere una distribuzione equilibrata degli oggetti sulla griglia. Ad esempio, potremmo posizionare gli oggetti 'I' nella posizione (2,5) e gli oggetti 'W' nella posizione (4,5). Questo ci darebbe un piano completo per tutti gli oggetti e potremmo iniziare a pianificare le mosse necessarie per raggiungere le posizioni finali. Pensiamo di procedere con questo piano e di iniziare a spostare gli oggetti secondo l'ordine che abbiamo discusso, partendo dagli oggetti 'H'?} \\
            }
        }
    }
    & & \\ \\

    \theutterance \stepcounter{utterance}  
    & & & \multicolumn{2}{p{0.3\linewidth}}{
        \cellcolor[rgb]{0.9,0.9,0.9}{
            \makecell[{{p{\linewidth}}}]{
                \texttt{\tiny{[GM$|$GM]}}
                \texttt{COMUNICA: Sono d'accordo con la posizione finale per gli oggetti 'H' e 'T' che hai proposto, (3,1) per 'H' e (3,7) per 'T'. Per gli oggetti 'I' e 'W', potremmo considerare di posizionarli nelle colonne 2 e 4, in modo da avere una distribuzione equilibrata degli oggetti sulla griglia. Ad esempio, potremmo posizionare gli oggetti 'I' nella posizione (2,5) e gli oggetti 'W' nella posizione (4,5). Questo ci darebbe un piano completo per tutti gli oggetti e potremmo iniziare a pianificare le mosse necessarie per raggiungere le posizioni finali. Pensiamo di procedere con questo piano e di iniziare a spostare gli oggetti secondo l'ordine che abbiamo discusso, partendo dagli oggetti 'H'?} \\
            }
        }
    }
    & & \\ \\

    \theutterance \stepcounter{utterance}  
    & & \multicolumn{4}{p{0.6\linewidth}}{
        \cellcolor[rgb]{0.9,0.9,0.9}{
            \makecell[{{p{\linewidth}}}]{
                \texttt{\tiny{[GM$\rangle$P2]}}
                \texttt{Il tuo messaggio è stato inoltrato all'altro giocatore.} \\
\\ 
\\ 
\texttt{You are currently playing round 5 of maximum 20.} \\
\texttt{You have collectively accumulated 2 of 12 penalties.} \\
\texttt{L'altro giocatore ha inviato questo messaggio:} \\
\\ 
\\ 
\texttt{"Sono d'accordo con la posizione finale per gli oggetti 'H' e 'T' che hai proposto, (3,1) per 'H' e (3,7) per 'T'. Per gli oggetti 'I' e 'W', potremmo considerare di posizionarli nelle colonne 2 e 4, in modo da avere una distribuzione equilibrata degli oggetti sulla griglia. Ad esempio, potremmo posizionare gli oggetti 'I' nella posizione (2,5) e gli oggetti 'W' nella posizione (4,5). Questo ci darebbe un piano completo per tutti gli oggetti e potremmo iniziare a pianificare le mosse necessarie per raggiungere le posizioni finali. Pensiamo di procedere con questo piano e di iniziare a spostare gli oggetti secondo l'ordine che abbiamo discusso, partendo dagli oggetti 'H'?"} \\
\\ 
\\ 
\texttt{Qual è il tuo prossimo comando?} \\
            }
        }
    }
    & \\ \\

    \theutterance \stepcounter{utterance}  
    & & & \multicolumn{4}{p{0.6\linewidth}}{
        \cellcolor[rgb]{0.9,0.9,0.9}{
            \makecell[{{p{\linewidth}}}]{
                \texttt{\tiny{[GM$\langle$P2]}}
                \texttt{COMUNICA: Sono d'accordo con il piano che hai proposto per gli oggetti 'I' e 'W', posizionarli nelle colonne 2 e 4 sembra una buona idea. La posizione (2,5) per gli oggetti 'I' e (4,5) per gli oggetti 'W' sembra ragionevole. Adesso che abbiamo un piano completo, possiamo iniziare a spostare gli oggetti secondo l'ordine discusso, partendo dagli oggetti 'H'. Per iniziare, potrei spostare il mio 'T' verso la posizione (3,7) come pianificato, per allinearli con il tuo 'T'. Come procedi con il tuo 'H' e gli altri oggetti?} \\
            }
        }
    }
     \\ \\

    \theutterance \stepcounter{utterance}  
    & & & \multicolumn{2}{p{0.3\linewidth}}{
        \cellcolor[rgb]{0.9,0.9,0.9}{
            \makecell[{{p{\linewidth}}}]{
                \texttt{\tiny{[GM$|$GM]}}
                \texttt{COMUNICA: Sono d'accordo con il piano che hai proposto per gli oggetti 'I' e 'W', posizionarli nelle colonne 2 e 4 sembra una buona idea. La posizione (2,5) per gli oggetti 'I' e (4,5) per gli oggetti 'W' sembra ragionevole. Adesso che abbiamo un piano completo, possiamo iniziare a spostare gli oggetti secondo l'ordine discusso, partendo dagli oggetti 'H'. Per iniziare, potrei spostare il mio 'T' verso la posizione (3,7) come pianificato, per allinearli con il tuo 'T'. Come procedi con il tuo 'H' e gli altri oggetti?} \\
            }
        }
    }
    & & \\ \\

    \theutterance \stepcounter{utterance}  
    & & \multicolumn{4}{p{0.6\linewidth}}{
        \cellcolor[rgb]{0.9,0.9,0.9}{
            \makecell[{{p{\linewidth}}}]{
                \texttt{\tiny{[P1$\langle$GM]}}
                \texttt{Il tuo messaggio è stato inoltrato all'altro giocatore.} \\
\\ 
\\ 
\texttt{You are currently playing round 5 of maximum 20.} \\
\texttt{You have collectively accumulated 2 of 12 penalties.} \\
\texttt{L'altro giocatore ha inviato questo messaggio:} \\
\\ 
\\ 
\texttt{"Sono d'accordo con il piano che hai proposto per gli oggetti 'I' e 'W', posizionarli nelle colonne 2 e 4 sembra una buona idea. La posizione (2,5) per gli oggetti 'I' e (4,5) per gli oggetti 'W' sembra ragionevole. Adesso che abbiamo un piano completo, possiamo iniziare a spostare gli oggetti secondo l'ordine discusso, partendo dagli oggetti 'H'. Per iniziare, potrei spostare il mio 'T' verso la posizione (3,7) come pianificato, per allinearli con il tuo 'T'. Come procedi con il tuo 'H' e gli altri oggetti?"} \\
\\ 
\\ 
\texttt{Qual è il tuo prossimo comando?} \\
            }
        }
    }
    & \\ \\

    \theutterance \stepcounter{utterance}  
    & \multicolumn{4}{p{0.6\linewidth}}{
        \cellcolor[rgb]{0.9,0.9,0.9}{
            \makecell[{{p{\linewidth}}}]{
                \texttt{\tiny{[P1$\rangle$GM]}}
                \texttt{SPOSTA: H, (3,6)} \\
            }
        }
    }
    & & \\ \\

    \theutterance \stepcounter{utterance}  
    & & & \multicolumn{2}{p{0.3\linewidth}}{
        \cellcolor[rgb]{0.9,0.9,0.9}{
            \makecell[{{p{\linewidth}}}]{
                \texttt{\tiny{[GM$|$GM]}}
                \texttt{SPOSTA: H, (3,6)} \\
            }
        }
    }
    & & \\ \\

    \theutterance \stepcounter{utterance}  
    & & & \multicolumn{2}{p{0.3\linewidth}}{
        \cellcolor[rgb]{0.9,0.9,0.9}{
            \makecell[{{p{\linewidth}}}]{
                \texttt{\tiny{[GM$|$GM]}}
                \texttt{Penalità: (3,6) non è vuoto, ma contiene '─'.} \\
\\ 
\texttt{You have collectively accumulated 3 of 12 penalties.} \\
\texttt{Riprova!} \\
            }
        }
    }
    & & \\ \\

    \theutterance \stepcounter{utterance}  
    & & \multicolumn{4}{p{0.6\linewidth}}{
        \cellcolor[rgb]{0.9,0.9,0.9}{
            \makecell[{{p{\linewidth}}}]{
                \texttt{\tiny{[P1$\langle$GM]}}
                \texttt{Penalità: (3,6) non è vuoto, ma contiene '─'.} \\
\\ 
\texttt{You have collectively accumulated 3 of 12 penalties.} \\
\texttt{Riprova!} \\
            }
        }
    }
    & \\ \\

    \theutterance \stepcounter{utterance}  
    & \multicolumn{4}{p{0.6\linewidth}}{
        \cellcolor[rgb]{0.9,0.9,0.9}{
            \makecell[{{p{\linewidth}}}]{
                \texttt{\tiny{[P1$\rangle$GM]}}
                \texttt{SPOSTA: H, (3,1)} \\
            }
        }
    }
    & & \\ \\

    \theutterance \stepcounter{utterance}  
    & & & \multicolumn{2}{p{0.3\linewidth}}{
        \cellcolor[rgb]{0.9,0.9,0.9}{
            \makecell[{{p{\linewidth}}}]{
                \texttt{\tiny{[GM$|$GM]}}
                \texttt{SPOSTA: H, (3,1)} \\
            }
        }
    }
    & & \\ \\

    \theutterance \stepcounter{utterance}  
    & & & \multicolumn{2}{p{0.3\linewidth}}{
        \cellcolor[rgb]{0.9,0.9,0.9}{
            \makecell[{{p{\linewidth}}}]{
                \texttt{\tiny{[GM$|$GM]}}
                \texttt{Spostato 'H' in (3,1) con successo. La tua griglia aggiornata appare così:} \\
\texttt{```} \\
\texttt{ 1234567} \\
\texttt{╔═══════╗} \\
\texttt{║I◌H◌◌◌◌║ 1} \\
\texttt{║◌◌◌◌◌◌◌║ 2} \\
\texttt{║◌◌◌T◌┌─╢ 3} \\
\texttt{║◌◌◌┌─┤◌║ 4} \\
\texttt{║◌◌W│C├─╢ 5} \\
\texttt{╟───┼─┘◌║ 6} \\
\texttt{║◌◌◌│◌◌◌║ 7} \\
\texttt{╚═══╧═══╝} \\
\texttt{```} \\
            }
        }
    }
    & & \\ \\

    \theutterance \stepcounter{utterance}  
    & & \multicolumn{4}{p{0.6\linewidth}}{
        \cellcolor[rgb]{0.9,0.9,0.9}{
            \makecell[{{p{\linewidth}}}]{
                \texttt{\tiny{[GM$\rangle$P2]}}
                \texttt{Il tuo messaggio è stato inoltrato all'altro giocatore.} \\
\\ 
\\ 
\texttt{You are currently playing round 6 of maximum 20.} \\
\texttt{You have collectively accumulated 3 of 12 penalties.} \\
\texttt{L'altro giocatore ha spostato un oggetto sulla sua griglia.} \\
\\ 
\\ 
\texttt{Qual è il tuo prossimo comando?} \\
            }
        }
    }
    & \\ \\

    \theutterance \stepcounter{utterance}  
    & & & \multicolumn{4}{p{0.6\linewidth}}{
        \cellcolor[rgb]{0.9,0.9,0.9}{
            \makecell[{{p{\linewidth}}}]{
                \texttt{\tiny{[GM$\langle$P2]}}
                \texttt{SPOSTA: T, (3,7)} \\
            }
        }
    }
     \\ \\

    \theutterance \stepcounter{utterance}  
    & & & \multicolumn{2}{p{0.3\linewidth}}{
        \cellcolor[rgb]{0.9,0.9,0.9}{
            \makecell[{{p{\linewidth}}}]{
                \texttt{\tiny{[GM$|$GM]}}
                \texttt{SPOSTA: T, (3,7)} \\
            }
        }
    }
    & & \\ \\

    \theutterance \stepcounter{utterance}  
    & & & \multicolumn{2}{p{0.3\linewidth}}{
        \cellcolor[rgb]{0.9,0.9,0.9}{
            \makecell[{{p{\linewidth}}}]{
                \texttt{\tiny{[GM$|$GM]}}
                \texttt{Spostato 'T' in (3,7) con successo. La tua griglia aggiornata appare così:} \\
\texttt{```} \\
\texttt{ 1234567} \\
\texttt{╔═══════╗} \\
\texttt{║◌◌H◌◌◌◌║ 1} \\
\texttt{║◌◌◌C◌◌◌║ 2} \\
\texttt{║◌◌◌◌◌┌─╢ 3} \\
\texttt{║◌◌◌┌─┤◌║ 4} \\
\texttt{║◌IW│◌├─╢ 5} \\
\texttt{╟───┼─┘◌║ 6} \\
\texttt{║◌◌T│◌◌◌║ 7} \\
\texttt{╚═══╧═══╝} \\
\texttt{```} \\
            }
        }
    }
    & & \\ \\

    \theutterance \stepcounter{utterance}  
    & & \multicolumn{4}{p{0.6\linewidth}}{
        \cellcolor[rgb]{0.9,0.9,0.9}{
            \makecell[{{p{\linewidth}}}]{
                \texttt{\tiny{[P1$\langle$GM]}}
                \texttt{Spostato 'H' in (3,1) con successo. La tua griglia aggiornata appare così:} \\
\texttt{```} \\
\texttt{ 1234567} \\
\texttt{╔═══════╗} \\
\texttt{║I◌H◌◌◌◌║ 1} \\
\texttt{║◌◌◌◌◌◌◌║ 2} \\
\texttt{║◌◌◌T◌┌─╢ 3} \\
\texttt{║◌◌◌┌─┤◌║ 4} \\
\texttt{║◌◌W│C├─╢ 5} \\
\texttt{╟───┼─┘◌║ 6} \\
\texttt{║◌◌◌│◌◌◌║ 7} \\
\texttt{╚═══╧═══╝} \\
\texttt{```} \\
\texttt{You are currently playing round 6 of maximum 20.} \\
\texttt{You have collectively accumulated 3 of 12 penalties.} \\
\texttt{L'altro giocatore ha spostato un oggetto sulla sua griglia.} \\
\\ 
\\ 
\texttt{Qual è il tuo prossimo comando?} \\
            }
        }
    }
    & \\ \\

    \theutterance \stepcounter{utterance}  
    & \multicolumn{4}{p{0.6\linewidth}}{
        \cellcolor[rgb]{0.9,0.9,0.9}{
            \makecell[{{p{\linewidth}}}]{
                \texttt{\tiny{[P1$\rangle$GM]}}
                \texttt{SPOSTA: T, (3,7)} \\
            }
        }
    }
    & & \\ \\

    \theutterance \stepcounter{utterance}  
    & & & \multicolumn{2}{p{0.3\linewidth}}{
        \cellcolor[rgb]{0.9,0.9,0.9}{
            \makecell[{{p{\linewidth}}}]{
                \texttt{\tiny{[GM$|$GM]}}
                \texttt{SPOSTA: T, (3,7)} \\
            }
        }
    }
    & & \\ \\

    \theutterance \stepcounter{utterance}  
    & & & \multicolumn{2}{p{0.3\linewidth}}{
        \cellcolor[rgb]{0.9,0.9,0.9}{
            \makecell[{{p{\linewidth}}}]{
                \texttt{\tiny{[GM$|$GM]}}
                \texttt{Spostato 'T' in (3,7) con successo. La tua griglia aggiornata appare così:} \\
\texttt{```} \\
\texttt{ 1234567} \\
\texttt{╔═══════╗} \\
\texttt{║I◌H◌◌◌◌║ 1} \\
\texttt{║◌◌◌◌◌◌◌║ 2} \\
\texttt{║◌◌◌◌◌┌─╢ 3} \\
\texttt{║◌◌◌┌─┤◌║ 4} \\
\texttt{║◌◌W│C├─╢ 5} \\
\texttt{╟───┼─┘◌║ 6} \\
\texttt{║◌◌T│◌◌◌║ 7} \\
\texttt{╚═══╧═══╝} \\
\texttt{```} \\
            }
        }
    }
    & & \\ \\

    \theutterance \stepcounter{utterance}  
    & & \multicolumn{4}{p{0.6\linewidth}}{
        \cellcolor[rgb]{0.9,0.9,0.9}{
            \makecell[{{p{\linewidth}}}]{
                \texttt{\tiny{[GM$\rangle$P2]}}
                \texttt{Spostato 'T' in (3,7) con successo. La tua griglia aggiornata appare così:} \\
\texttt{```} \\
\texttt{ 1234567} \\
\texttt{╔═══════╗} \\
\texttt{║◌◌H◌◌◌◌║ 1} \\
\texttt{║◌◌◌C◌◌◌║ 2} \\
\texttt{║◌◌◌◌◌┌─╢ 3} \\
\texttt{║◌◌◌┌─┤◌║ 4} \\
\texttt{║◌IW│◌├─╢ 5} \\
\texttt{╟───┼─┘◌║ 6} \\
\texttt{║◌◌T│◌◌◌║ 7} \\
\texttt{╚═══╧═══╝} \\
\texttt{```} \\
\texttt{You are currently playing round 7 of maximum 20.} \\
\texttt{You have collectively accumulated 3 of 12 penalties.} \\
\texttt{L'altro giocatore ha spostato un oggetto sulla sua griglia.} \\
\\ 
\\ 
\texttt{Qual è il tuo prossimo comando?} \\
            }
        }
    }
    & \\ \\

    \theutterance \stepcounter{utterance}  
    & & & \multicolumn{4}{p{0.6\linewidth}}{
        \cellcolor[rgb]{0.9,0.9,0.9}{
            \makecell[{{p{\linewidth}}}]{
                \texttt{\tiny{[GM$\langle$P2]}}
                \texttt{COMUNICA: Sembra che entrambi abbiamo spostato i nostri oggetti 'T' verso la posizione desiderata. Ora che gli oggetti 'T' sono allineati, possiamo procedere con gli oggetti 'H'. Come stanno gli oggetti 'H' sulla tua griglia? Hai spostato il tuo 'H' nella posizione (3,1) come discusso?} \\
            }
        }
    }
     \\ \\

    \theutterance \stepcounter{utterance}  
    & & & \multicolumn{2}{p{0.3\linewidth}}{
        \cellcolor[rgb]{0.9,0.9,0.9}{
            \makecell[{{p{\linewidth}}}]{
                \texttt{\tiny{[GM$|$GM]}}
                \texttt{COMUNICA: Sembra che entrambi abbiamo spostato i nostri oggetti 'T' verso la posizione desiderata. Ora che gli oggetti 'T' sono allineati, possiamo procedere con gli oggetti 'H'. Come stanno gli oggetti 'H' sulla tua griglia? Hai spostato il tuo 'H' nella posizione (3,1) come discusso?} \\
            }
        }
    }
    & & \\ \\

    \theutterance \stepcounter{utterance}  
    & & \multicolumn{4}{p{0.6\linewidth}}{
        \cellcolor[rgb]{0.9,0.9,0.9}{
            \makecell[{{p{\linewidth}}}]{
                \texttt{\tiny{[P1$\langle$GM]}}
                \texttt{Spostato 'T' in (3,7) con successo. La tua griglia aggiornata appare così:} \\
\texttt{```} \\
\texttt{ 1234567} \\
\texttt{╔═══════╗} \\
\texttt{║I◌H◌◌◌◌║ 1} \\
\texttt{║◌◌◌◌◌◌◌║ 2} \\
\texttt{║◌◌◌◌◌┌─╢ 3} \\
\texttt{║◌◌◌┌─┤◌║ 4} \\
\texttt{║◌◌W│C├─╢ 5} \\
\texttt{╟───┼─┘◌║ 6} \\
\texttt{║◌◌T│◌◌◌║ 7} \\
\texttt{╚═══╧═══╝} \\
\texttt{```} \\
\texttt{You are currently playing round 7 of maximum 20.} \\
\texttt{You have collectively accumulated 3 of 12 penalties.} \\
\texttt{L'altro giocatore ha inviato questo messaggio:} \\
\\ 
\\ 
\texttt{"Sembra che entrambi abbiamo spostato i nostri oggetti 'T' verso la posizione desiderata. Ora che gli oggetti 'T' sono allineati, possiamo procedere con gli oggetti 'H'. Come stanno gli oggetti 'H' sulla tua griglia? Hai spostato il tuo 'H' nella posizione (3,1) come discusso?"} \\
\\ 
\\ 
\texttt{Qual è il tuo prossimo comando?} \\
            }
        }
    }
    & \\ \\

    \theutterance \stepcounter{utterance}  
    & \multicolumn{4}{p{0.6\linewidth}}{
        \cellcolor[rgb]{0.9,0.9,0.9}{
            \makecell[{{p{\linewidth}}}]{
                \texttt{\tiny{[P1$\rangle$GM]}}
                \texttt{COMUNICA: Sì, ho spostato il mio 'H' nella posizione (3,1) come pianificato. Ora che gli oggetti 'T' sono allineati, possiamo procedere con gli oggetti 'I' e 'W'. Sono pronto a spostare i miei 'I' e 'W' nelle posizioni che abbiamo discusso, (2,5) e (4,5) rispettivamente. Come procedi con i tuoi oggetti 'I' e 'W'?} \\
            }
        }
    }
    & & \\ \\

    \theutterance \stepcounter{utterance}  
    & & & \multicolumn{2}{p{0.3\linewidth}}{
        \cellcolor[rgb]{0.9,0.9,0.9}{
            \makecell[{{p{\linewidth}}}]{
                \texttt{\tiny{[GM$|$GM]}}
                \texttt{COMUNICA: Sì, ho spostato il mio 'H' nella posizione (3,1) come pianificato. Ora che gli oggetti 'T' sono allineati, possiamo procedere con gli oggetti 'I' e 'W'. Sono pronto a spostare i miei 'I' e 'W' nelle posizioni che abbiamo discusso, (2,5) e (4,5) rispettivamente. Come procedi con i tuoi oggetti 'I' e 'W'?} \\
            }
        }
    }
    & & \\ \\

    \theutterance \stepcounter{utterance}  
    & & \multicolumn{4}{p{0.6\linewidth}}{
        \cellcolor[rgb]{0.9,0.9,0.9}{
            \makecell[{{p{\linewidth}}}]{
                \texttt{\tiny{[GM$\rangle$P2]}}
                \texttt{Il tuo messaggio è stato inoltrato all'altro giocatore.} \\
\\ 
\\ 
\texttt{You are currently playing round 8 of maximum 20.} \\
\texttt{You have collectively accumulated 3 of 12 penalties.} \\
\texttt{L'altro giocatore ha inviato questo messaggio:} \\
\\ 
\\ 
\texttt{"Sì, ho spostato il mio 'H' nella posizione (3,1) come pianificato. Ora che gli oggetti 'T' sono allineati, possiamo procedere con gli oggetti 'I' e 'W'. Sono pronto a spostare i miei 'I' e 'W' nelle posizioni che abbiamo discusso, (2,5) e (4,5) rispettivamente. Come procedi con i tuoi oggetti 'I' e 'W'?"} \\
\\ 
\\ 
\texttt{Qual è il tuo prossimo comando?} \\
            }
        }
    }
    & \\ \\

    \theutterance \stepcounter{utterance}  
    & & & \multicolumn{4}{p{0.6\linewidth}}{
        \cellcolor[rgb]{0.9,0.9,0.9}{
            \makecell[{{p{\linewidth}}}]{
                \texttt{\tiny{[GM$\langle$P2]}}
                \texttt{SPOSTA: I, (2,5)} \\
            }
        }
    }
     \\ \\

    \theutterance \stepcounter{utterance}  
    & & & \multicolumn{2}{p{0.3\linewidth}}{
        \cellcolor[rgb]{0.9,0.9,0.9}{
            \makecell[{{p{\linewidth}}}]{
                \texttt{\tiny{[GM$|$GM]}}
                \texttt{SPOSTA: I, (2,5)} \\
            }
        }
    }
    & & \\ \\

    \theutterance \stepcounter{utterance}  
    & & & \multicolumn{2}{p{0.3\linewidth}}{
        \cellcolor[rgb]{0.9,0.9,0.9}{
            \makecell[{{p{\linewidth}}}]{
                \texttt{\tiny{[GM$|$GM]}}
                \texttt{Penalità: (2,5) non è vuoto, ma contiene 'I'.} \\
\\ 
\texttt{You have collectively accumulated 4 of 12 penalties.} \\
\texttt{Riprova!} \\
            }
        }
    }
    & & \\ \\

    \theutterance \stepcounter{utterance}  
    & & \multicolumn{4}{p{0.6\linewidth}}{
        \cellcolor[rgb]{0.9,0.9,0.9}{
            \makecell[{{p{\linewidth}}}]{
                \texttt{\tiny{[GM$\rangle$P2]}}
                \texttt{Penalità: (2,5) non è vuoto, ma contiene 'I'.} \\
\\ 
\texttt{You have collectively accumulated 4 of 12 penalties.} \\
\texttt{Riprova!} \\
            }
        }
    }
    & \\ \\

    \theutterance \stepcounter{utterance}  
    & & & \multicolumn{4}{p{0.6\linewidth}}{
        \cellcolor[rgb]{0.9,0.9,0.9}{
            \makecell[{{p{\linewidth}}}]{
                \texttt{\tiny{[GM$\langle$P2]}}
                \texttt{SPOSTA: W, (4,5)} \\
            }
        }
    }
     \\ \\

    \theutterance \stepcounter{utterance}  
    & & & \multicolumn{2}{p{0.3\linewidth}}{
        \cellcolor[rgb]{0.9,0.9,0.9}{
            \makecell[{{p{\linewidth}}}]{
                \texttt{\tiny{[GM$|$GM]}}
                \texttt{SPOSTA: W, (4,5)} \\
            }
        }
    }
    & & \\ \\

    \theutterance \stepcounter{utterance}  
    & & & \multicolumn{2}{p{0.3\linewidth}}{
        \cellcolor[rgb]{0.9,0.9,0.9}{
            \makecell[{{p{\linewidth}}}]{
                \texttt{\tiny{[GM$|$GM]}}
                \texttt{Penalità: (4,5) non è vuoto, ma contiene '│'.} \\
\\ 
\texttt{You have collectively accumulated 5 of 12 penalties.} \\
\texttt{Riprova!} \\
            }
        }
    }
    & & \\ \\

    \theutterance \stepcounter{utterance}  
    & & \multicolumn{4}{p{0.6\linewidth}}{
        \cellcolor[rgb]{0.9,0.9,0.9}{
            \makecell[{{p{\linewidth}}}]{
                \texttt{\tiny{[GM$\rangle$P2]}}
                \texttt{Penalità: (4,5) non è vuoto, ma contiene '│'.} \\
\\ 
\texttt{You have collectively accumulated 5 of 12 penalties.} \\
\texttt{Riprova!} \\
            }
        }
    }
    & \\ \\

    \theutterance \stepcounter{utterance}  
    & & & \multicolumn{4}{p{0.6\linewidth}}{
        \cellcolor[rgb]{0.9,0.9,0.9}{
            \makecell[{{p{\linewidth}}}]{
                \texttt{\tiny{[GM$\langle$P2]}}
                \texttt{SPOSTA: I, (2,4)} \\
            }
        }
    }
     \\ \\

    \theutterance \stepcounter{utterance}  
    & & & \multicolumn{2}{p{0.3\linewidth}}{
        \cellcolor[rgb]{0.9,0.9,0.9}{
            \makecell[{{p{\linewidth}}}]{
                \texttt{\tiny{[GM$|$GM]}}
                \texttt{SPOSTA: I, (2,4)} \\
            }
        }
    }
    & & \\ \\

    \theutterance \stepcounter{utterance}  
    & & & \multicolumn{2}{p{0.3\linewidth}}{
        \cellcolor[rgb]{0.9,0.9,0.9}{
            \makecell[{{p{\linewidth}}}]{
                \texttt{\tiny{[GM$|$GM]}}
                \texttt{Spostato 'I' in (2,4) con successo. La tua griglia aggiornata appare così:} \\
\texttt{```} \\
\texttt{ 1234567} \\
\texttt{╔═══════╗} \\
\texttt{║◌◌H◌◌◌◌║ 1} \\
\texttt{║◌◌◌C◌◌◌║ 2} \\
\texttt{║◌◌◌◌◌┌─╢ 3} \\
\texttt{║◌I◌┌─┤◌║ 4} \\
\texttt{║◌◌W│◌├─╢ 5} \\
\texttt{╟───┼─┘◌║ 6} \\
\texttt{║◌◌T│◌◌◌║ 7} \\
\texttt{╚═══╧═══╝} \\
\texttt{```} \\
            }
        }
    }
    & & \\ \\

    \theutterance \stepcounter{utterance}  
    & & \multicolumn{4}{p{0.6\linewidth}}{
        \cellcolor[rgb]{0.9,0.9,0.9}{
            \makecell[{{p{\linewidth}}}]{
                \texttt{\tiny{[P1$\langle$GM]}}
                \texttt{Il tuo messaggio è stato inoltrato all'altro giocatore.} \\
\\ 
\\ 
\texttt{You are currently playing round 8 of maximum 20.} \\
\texttt{You have collectively accumulated 5 of 12 penalties.} \\
\texttt{L'altro giocatore ha spostato un oggetto sulla sua griglia.} \\
\\ 
\\ 
\texttt{Qual è il tuo prossimo comando?} \\
            }
        }
    }
    & \\ \\

    \theutterance \stepcounter{utterance}  
    & \multicolumn{4}{p{0.6\linewidth}}{
        \cellcolor[rgb]{0.9,0.9,0.9}{
            \makecell[{{p{\linewidth}}}]{
                \texttt{\tiny{[P1$\rangle$GM]}}
                \texttt{SPOSTA: I, (2,5)} \\
            }
        }
    }
    & & \\ \\

    \theutterance \stepcounter{utterance}  
    & & & \multicolumn{2}{p{0.3\linewidth}}{
        \cellcolor[rgb]{0.9,0.9,0.9}{
            \makecell[{{p{\linewidth}}}]{
                \texttt{\tiny{[GM$|$GM]}}
                \texttt{SPOSTA: I, (2,5)} \\
            }
        }
    }
    & & \\ \\

    \theutterance \stepcounter{utterance}  
    & & & \multicolumn{2}{p{0.3\linewidth}}{
        \cellcolor[rgb]{0.9,0.9,0.9}{
            \makecell[{{p{\linewidth}}}]{
                \texttt{\tiny{[GM$|$GM]}}
                \texttt{Spostato 'I' in (2,5) con successo. La tua griglia aggiornata appare così:} \\
\texttt{```} \\
\texttt{ 1234567} \\
\texttt{╔═══════╗} \\
\texttt{║◌◌H◌◌◌◌║ 1} \\
\texttt{║◌◌◌◌◌◌◌║ 2} \\
\texttt{║◌◌◌◌◌┌─╢ 3} \\
\texttt{║◌◌◌┌─┤◌║ 4} \\
\texttt{║◌IW│C├─╢ 5} \\
\texttt{╟───┼─┘◌║ 6} \\
\texttt{║◌◌T│◌◌◌║ 7} \\
\texttt{╚═══╧═══╝} \\
\texttt{```} \\
            }
        }
    }
    & & \\ \\

    \theutterance \stepcounter{utterance}  
    & & \multicolumn{4}{p{0.6\linewidth}}{
        \cellcolor[rgb]{0.9,0.9,0.9}{
            \makecell[{{p{\linewidth}}}]{
                \texttt{\tiny{[GM$\rangle$P2]}}
                \texttt{Spostato 'I' in (2,4) con successo. La tua griglia aggiornata appare così:} \\
\texttt{```} \\
\texttt{ 1234567} \\
\texttt{╔═══════╗} \\
\texttt{║◌◌H◌◌◌◌║ 1} \\
\texttt{║◌◌◌C◌◌◌║ 2} \\
\texttt{║◌◌◌◌◌┌─╢ 3} \\
\texttt{║◌I◌┌─┤◌║ 4} \\
\texttt{║◌◌W│◌├─╢ 5} \\
\texttt{╟───┼─┘◌║ 6} \\
\texttt{║◌◌T│◌◌◌║ 7} \\
\texttt{╚═══╧═══╝} \\
\texttt{```} \\
\texttt{You are currently playing round 9 of maximum 20.} \\
\texttt{You have collectively accumulated 5 of 12 penalties.} \\
\texttt{L'altro giocatore ha spostato un oggetto sulla sua griglia.} \\
\\ 
\\ 
\texttt{Qual è il tuo prossimo comando?} \\
            }
        }
    }
    & \\ \\

    \theutterance \stepcounter{utterance}  
    & & & \multicolumn{4}{p{0.6\linewidth}}{
        \cellcolor[rgb]{0.9,0.9,0.9}{
            \makecell[{{p{\linewidth}}}]{
                \texttt{\tiny{[GM$\langle$P2]}}
                \texttt{SPOSTA: W, (4,5)} \\
            }
        }
    }
     \\ \\

    \theutterance \stepcounter{utterance}  
    & & & \multicolumn{2}{p{0.3\linewidth}}{
        \cellcolor[rgb]{0.9,0.9,0.9}{
            \makecell[{{p{\linewidth}}}]{
                \texttt{\tiny{[GM$|$GM]}}
                \texttt{SPOSTA: W, (4,5)} \\
            }
        }
    }
    & & \\ \\

    \theutterance \stepcounter{utterance}  
    & & & \multicolumn{2}{p{0.3\linewidth}}{
        \cellcolor[rgb]{0.9,0.9,0.9}{
            \makecell[{{p{\linewidth}}}]{
                \texttt{\tiny{[GM$|$GM]}}
                \texttt{Penalità: (4,5) non è vuoto, ma contiene '│'.} \\
\\ 
\texttt{You have collectively accumulated 6 of 12 penalties.} \\
\texttt{Riprova!} \\
            }
        }
    }
    & & \\ \\

    \theutterance \stepcounter{utterance}  
    & & \multicolumn{4}{p{0.6\linewidth}}{
        \cellcolor[rgb]{0.9,0.9,0.9}{
            \makecell[{{p{\linewidth}}}]{
                \texttt{\tiny{[GM$\rangle$P2]}}
                \texttt{Penalità: (4,5) non è vuoto, ma contiene '│'.} \\
\\ 
\texttt{You have collectively accumulated 6 of 12 penalties.} \\
\texttt{Riprova!} \\
            }
        }
    }
    & \\ \\

    \theutterance \stepcounter{utterance}  
    & & & \multicolumn{4}{p{0.6\linewidth}}{
        \cellcolor[rgb]{0.9,0.9,0.9}{
            \makecell[{{p{\linewidth}}}]{
                \texttt{\tiny{[GM$\langle$P2]}}
                \texttt{SPOSTA: W, (4,4)} \\
            }
        }
    }
     \\ \\

    \theutterance \stepcounter{utterance}  
    & & & \multicolumn{2}{p{0.3\linewidth}}{
        \cellcolor[rgb]{0.9,0.9,0.9}{
            \makecell[{{p{\linewidth}}}]{
                \texttt{\tiny{[GM$|$GM]}}
                \texttt{SPOSTA: W, (4,4)} \\
            }
        }
    }
    & & \\ \\

    \theutterance \stepcounter{utterance}  
    & & & \multicolumn{2}{p{0.3\linewidth}}{
        \cellcolor[rgb]{0.9,0.9,0.9}{
            \makecell[{{p{\linewidth}}}]{
                \texttt{\tiny{[GM$|$GM]}}
                \texttt{Penalità: (4,4) non è vuoto, ma contiene '┌'.} \\
\\ 
\texttt{You have collectively accumulated 7 of 12 penalties.} \\
\texttt{Riprova!} \\
            }
        }
    }
    & & \\ \\

    \theutterance \stepcounter{utterance}  
    & & \multicolumn{4}{p{0.6\linewidth}}{
        \cellcolor[rgb]{0.9,0.9,0.9}{
            \makecell[{{p{\linewidth}}}]{
                \texttt{\tiny{[GM$\rangle$P2]}}
                \texttt{Penalità: (4,4) non è vuoto, ma contiene '┌'.} \\
\\ 
\texttt{You have collectively accumulated 7 of 12 penalties.} \\
\texttt{Riprova!} \\
            }
        }
    }
    & \\ \\

    \theutterance \stepcounter{utterance}  
    & & & \multicolumn{4}{p{0.6\linewidth}}{
        \cellcolor[rgb]{0.9,0.9,0.9}{
            \makecell[{{p{\linewidth}}}]{
                \texttt{\tiny{[GM$\langle$P2]}}
                \texttt{SPOSTA: W, (5,5)} \\
            }
        }
    }
     \\ \\

    \theutterance \stepcounter{utterance}  
    & & & \multicolumn{2}{p{0.3\linewidth}}{
        \cellcolor[rgb]{0.9,0.9,0.9}{
            \makecell[{{p{\linewidth}}}]{
                \texttt{\tiny{[GM$|$GM]}}
                \texttt{SPOSTA: W, (5,5)} \\
            }
        }
    }
    & & \\ \\

    \theutterance \stepcounter{utterance}  
    & & & \multicolumn{2}{p{0.3\linewidth}}{
        \cellcolor[rgb]{0.9,0.9,0.9}{
            \makecell[{{p{\linewidth}}}]{
                \texttt{\tiny{[GM$|$GM]}}
                \texttt{Spostato 'W' in (5,5) con successo. La tua griglia aggiornata appare così:} \\
\texttt{```} \\
\texttt{ 1234567} \\
\texttt{╔═══════╗} \\
\texttt{║◌◌H◌◌◌◌║ 1} \\
\texttt{║◌◌◌C◌◌◌║ 2} \\
\texttt{║◌◌◌◌◌┌─╢ 3} \\
\texttt{║◌I◌┌─┤◌║ 4} \\
\texttt{║◌◌◌│W├─╢ 5} \\
\texttt{╟───┼─┘◌║ 6} \\
\texttt{║◌◌T│◌◌◌║ 7} \\
\texttt{╚═══╧═══╝} \\
\texttt{```} \\
            }
        }
    }
    & & \\ \\

    \theutterance \stepcounter{utterance}  
    & & \multicolumn{4}{p{0.6\linewidth}}{
        \cellcolor[rgb]{0.9,0.9,0.9}{
            \makecell[{{p{\linewidth}}}]{
                \texttt{\tiny{[P1$\langle$GM]}}
                \texttt{Spostato 'I' in (2,5) con successo. La tua griglia aggiornata appare così:} \\
\texttt{```} \\
\texttt{ 1234567} \\
\texttt{╔═══════╗} \\
\texttt{║◌◌H◌◌◌◌║ 1} \\
\texttt{║◌◌◌◌◌◌◌║ 2} \\
\texttt{║◌◌◌◌◌┌─╢ 3} \\
\texttt{║◌◌◌┌─┤◌║ 4} \\
\texttt{║◌IW│C├─╢ 5} \\
\texttt{╟───┼─┘◌║ 6} \\
\texttt{║◌◌T│◌◌◌║ 7} \\
\texttt{╚═══╧═══╝} \\
\texttt{```} \\
\texttt{You are currently playing round 9 of maximum 20.} \\
\texttt{You have collectively accumulated 7 of 12 penalties.} \\
\texttt{L'altro giocatore ha spostato un oggetto sulla sua griglia.} \\
\\ 
\\ 
\texttt{Qual è il tuo prossimo comando?} \\
            }
        }
    }
    & \\ \\

    \theutterance \stepcounter{utterance}  
    & \multicolumn{4}{p{0.6\linewidth}}{
        \cellcolor[rgb]{0.9,0.9,0.9}{
            \makecell[{{p{\linewidth}}}]{
                \texttt{\tiny{[P1$\rangle$GM]}}
                \texttt{SPOSTA: W, (4,5)} \\
            }
        }
    }
    & & \\ \\

    \theutterance \stepcounter{utterance}  
    & & & \multicolumn{2}{p{0.3\linewidth}}{
        \cellcolor[rgb]{0.9,0.9,0.9}{
            \makecell[{{p{\linewidth}}}]{
                \texttt{\tiny{[GM$|$GM]}}
                \texttt{SPOSTA: W, (4,5)} \\
            }
        }
    }
    & & \\ \\

    \theutterance \stepcounter{utterance}  
    & & & \multicolumn{2}{p{0.3\linewidth}}{
        \cellcolor[rgb]{0.9,0.9,0.9}{
            \makecell[{{p{\linewidth}}}]{
                \texttt{\tiny{[GM$|$GM]}}
                \texttt{Penalità: (4,5) non è vuoto, ma contiene '│'.} \\
\\ 
\texttt{You have collectively accumulated 8 of 12 penalties.} \\
\texttt{Riprova!} \\
            }
        }
    }
    & & \\ \\

    \theutterance \stepcounter{utterance}  
    & & \multicolumn{4}{p{0.6\linewidth}}{
        \cellcolor[rgb]{0.9,0.9,0.9}{
            \makecell[{{p{\linewidth}}}]{
                \texttt{\tiny{[P1$\langle$GM]}}
                \texttt{Penalità: (4,5) non è vuoto, ma contiene '│'.} \\
\\ 
\texttt{You have collectively accumulated 8 of 12 penalties.} \\
\texttt{Riprova!} \\
            }
        }
    }
    & \\ \\

    \theutterance \stepcounter{utterance}  
    & \multicolumn{4}{p{0.6\linewidth}}{
        \cellcolor[rgb]{0.9,0.9,0.9}{
            \makecell[{{p{\linewidth}}}]{
                \texttt{\tiny{[P1$\rangle$GM]}}
                \texttt{SPOSTA: W, (4,4)} \\
            }
        }
    }
    & & \\ \\

    \theutterance \stepcounter{utterance}  
    & & & \multicolumn{2}{p{0.3\linewidth}}{
        \cellcolor[rgb]{0.9,0.9,0.9}{
            \makecell[{{p{\linewidth}}}]{
                \texttt{\tiny{[GM$|$GM]}}
                \texttt{SPOSTA: W, (4,4)} \\
            }
        }
    }
    & & \\ \\

    \theutterance \stepcounter{utterance}  
    & & & \multicolumn{2}{p{0.3\linewidth}}{
        \cellcolor[rgb]{0.9,0.9,0.9}{
            \makecell[{{p{\linewidth}}}]{
                \texttt{\tiny{[GM$|$GM]}}
                \texttt{Penalità: (4,4) non è vuoto, ma contiene '┌'.} \\
\\ 
\texttt{You have collectively accumulated 9 of 12 penalties.} \\
\texttt{Riprova!} \\
            }
        }
    }
    & & \\ \\

    \theutterance \stepcounter{utterance}  
    & & \multicolumn{4}{p{0.6\linewidth}}{
        \cellcolor[rgb]{0.9,0.9,0.9}{
            \makecell[{{p{\linewidth}}}]{
                \texttt{\tiny{[P1$\langle$GM]}}
                \texttt{Penalità: (4,4) non è vuoto, ma contiene '┌'.} \\
\\ 
\texttt{You have collectively accumulated 9 of 12 penalties.} \\
\texttt{Riprova!} \\
            }
        }
    }
    & \\ \\

    \theutterance \stepcounter{utterance}  
    & \multicolumn{4}{p{0.6\linewidth}}{
        \cellcolor[rgb]{0.9,0.9,0.9}{
            \makecell[{{p{\linewidth}}}]{
                \texttt{\tiny{[P1$\rangle$GM]}}
                \texttt{SPOSTA: W, (4,3)} \\
            }
        }
    }
    & & \\ \\

    \theutterance \stepcounter{utterance}  
    & & & \multicolumn{2}{p{0.3\linewidth}}{
        \cellcolor[rgb]{0.9,0.9,0.9}{
            \makecell[{{p{\linewidth}}}]{
                \texttt{\tiny{[GM$|$GM]}}
                \texttt{SPOSTA: W, (4,3)} \\
            }
        }
    }
    & & \\ \\

    \theutterance \stepcounter{utterance}  
    & & & \multicolumn{2}{p{0.3\linewidth}}{
        \cellcolor[rgb]{0.9,0.9,0.9}{
            \makecell[{{p{\linewidth}}}]{
                \texttt{\tiny{[GM$|$GM]}}
                \texttt{Spostato 'W' in (4,3) con successo. La tua griglia aggiornata appare così:} \\
\texttt{```} \\
\texttt{ 1234567} \\
\texttt{╔═══════╗} \\
\texttt{║◌◌H◌◌◌◌║ 1} \\
\texttt{║◌◌◌◌◌◌◌║ 2} \\
\texttt{║◌◌◌W◌┌─╢ 3} \\
\texttt{║◌◌◌┌─┤◌║ 4} \\
\texttt{║◌I◌│C├─╢ 5} \\
\texttt{╟───┼─┘◌║ 6} \\
\texttt{║◌◌T│◌◌◌║ 7} \\
\texttt{╚═══╧═══╝} \\
\texttt{```} \\
            }
        }
    }
    & & \\ \\

    \theutterance \stepcounter{utterance}  
    & & \multicolumn{4}{p{0.6\linewidth}}{
        \cellcolor[rgb]{0.9,0.9,0.9}{
            \makecell[{{p{\linewidth}}}]{
                \texttt{\tiny{[GM$\rangle$P2]}}
                \texttt{Spostato 'W' in (5,5) con successo. La tua griglia aggiornata appare così:} \\
\texttt{```} \\
\texttt{ 1234567} \\
\texttt{╔═══════╗} \\
\texttt{║◌◌H◌◌◌◌║ 1} \\
\texttt{║◌◌◌C◌◌◌║ 2} \\
\texttt{║◌◌◌◌◌┌─╢ 3} \\
\texttt{║◌I◌┌─┤◌║ 4} \\
\texttt{║◌◌◌│W├─╢ 5} \\
\texttt{╟───┼─┘◌║ 6} \\
\texttt{║◌◌T│◌◌◌║ 7} \\
\texttt{╚═══╧═══╝} \\
\texttt{```} \\
\texttt{You are currently playing round 10 of maximum 20.} \\
\texttt{You have collectively accumulated 9 of 12 penalties.} \\
\texttt{L'altro giocatore ha spostato un oggetto sulla sua griglia.} \\
\\ 
\\ 
\texttt{Qual è il tuo prossimo comando?} \\
            }
        }
    }
    & \\ \\

    \theutterance \stepcounter{utterance}  
    & & & \multicolumn{4}{p{0.6\linewidth}}{
        \cellcolor[rgb]{0.9,0.9,0.9}{
            \makecell[{{p{\linewidth}}}]{
                \texttt{\tiny{[GM$\langle$P2]}}
                \texttt{COMUNICA: Sembra che la griglia stia iniziando a prendere forma. Vorrei verificare lo stato attuale degli oggetti 'H' e assicurarmi che siano allineati correttamente. Come stanno gli oggetti 'H' sulla tua griglia? Sono ancora nella posizione (3,1) come pianificato?} \\
            }
        }
    }
     \\ \\

    \theutterance \stepcounter{utterance}  
    & & & \multicolumn{2}{p{0.3\linewidth}}{
        \cellcolor[rgb]{0.9,0.9,0.9}{
            \makecell[{{p{\linewidth}}}]{
                \texttt{\tiny{[GM$|$GM]}}
                \texttt{COMUNICA: Sembra che la griglia stia iniziando a prendere forma. Vorrei verificare lo stato attuale degli oggetti 'H' e assicurarmi che siano allineati correttamente. Come stanno gli oggetti 'H' sulla tua griglia? Sono ancora nella posizione (3,1) come pianificato?} \\
            }
        }
    }
    & & \\ \\

    \theutterance \stepcounter{utterance}  
    & & \multicolumn{4}{p{0.6\linewidth}}{
        \cellcolor[rgb]{0.9,0.9,0.9}{
            \makecell[{{p{\linewidth}}}]{
                \texttt{\tiny{[P1$\langle$GM]}}
                \texttt{Spostato 'W' in (4,3) con successo. La tua griglia aggiornata appare così:} \\
\texttt{```} \\
\texttt{ 1234567} \\
\texttt{╔═══════╗} \\
\texttt{║◌◌H◌◌◌◌║ 1} \\
\texttt{║◌◌◌◌◌◌◌║ 2} \\
\texttt{║◌◌◌W◌┌─╢ 3} \\
\texttt{║◌◌◌┌─┤◌║ 4} \\
\texttt{║◌I◌│C├─╢ 5} \\
\texttt{╟───┼─┘◌║ 6} \\
\texttt{║◌◌T│◌◌◌║ 7} \\
\texttt{╚═══╧═══╝} \\
\texttt{```} \\
\texttt{You are currently playing round 10 of maximum 20.} \\
\texttt{You have collectively accumulated 9 of 12 penalties.} \\
\texttt{L'altro giocatore ha inviato questo messaggio:} \\
\\ 
\\ 
\texttt{"Sembra che la griglia stia iniziando a prendere forma. Vorrei verificare lo stato attuale degli oggetti 'H' e assicurarmi che siano allineati correttamente. Come stanno gli oggetti 'H' sulla tua griglia? Sono ancora nella posizione (3,1) come pianificato?"} \\
\\ 
\\ 
\texttt{Qual è il tuo prossimo comando?} \\
            }
        }
    }
    & \\ \\

    \theutterance \stepcounter{utterance}  
    & \multicolumn{4}{p{0.6\linewidth}}{
        \cellcolor[rgb]{0.9,0.9,0.9}{
            \makecell[{{p{\linewidth}}}]{
                \texttt{\tiny{[P1$\rangle$GM]}}
                \texttt{COMUNICA: Sì, il mio 'H' è ancora nella posizione (3,1) come pianificato. La mia griglia attuale è aggiornata e gli oggetti 'H' sono allineati con quelli sulla tua griglia. Procediamo con gli oggetti 'I' e 'W' per assicurarci che siano nelle posizioni corrette. Il mio 'I' è già nella posizione (2,5) e il mio 'W' è stato spostato nella posizione (4,3). Come stanno gli oggetti 'I' e 'W' sulla tua griglia?} \\
            }
        }
    }
    & & \\ \\

    \theutterance \stepcounter{utterance}  
    & & & \multicolumn{2}{p{0.3\linewidth}}{
        \cellcolor[rgb]{0.9,0.9,0.9}{
            \makecell[{{p{\linewidth}}}]{
                \texttt{\tiny{[GM$|$GM]}}
                \texttt{COMUNICA: Sì, il mio 'H' è ancora nella posizione (3,1) come pianificato. La mia griglia attuale è aggiornata e gli oggetti 'H' sono allineati con quelli sulla tua griglia. Procediamo con gli oggetti 'I' e 'W' per assicurarci che siano nelle posizioni corrette. Il mio 'I' è già nella posizione (2,5) e il mio 'W' è stato spostato nella posizione (4,3). Come stanno gli oggetti 'I' e 'W' sulla tua griglia?} \\
            }
        }
    }
    & & \\ \\

    \theutterance \stepcounter{utterance}  
    & & \multicolumn{4}{p{0.6\linewidth}}{
        \cellcolor[rgb]{0.9,0.9,0.9}{
            \makecell[{{p{\linewidth}}}]{
                \texttt{\tiny{[GM$\rangle$P2]}}
                \texttt{Il tuo messaggio è stato inoltrato all'altro giocatore.} \\
\\ 
\\ 
\texttt{You are currently playing round 11 of maximum 20.} \\
\texttt{You have collectively accumulated 9 of 12 penalties.} \\
\texttt{L'altro giocatore ha inviato questo messaggio:} \\
\\ 
\\ 
\texttt{"Sì, il mio 'H' è ancora nella posizione (3,1) come pianificato. La mia griglia attuale è aggiornata e gli oggetti 'H' sono allineati con quelli sulla tua griglia. Procediamo con gli oggetti 'I' e 'W' per assicurarci che siano nelle posizioni corrette. Il mio 'I' è già nella posizione (2,5) e il mio 'W' è stato spostato nella posizione (4,3). Come stanno gli oggetti 'I' e 'W' sulla tua griglia?"} \\
\\ 
\\ 
\texttt{Qual è il tuo prossimo comando?} \\
            }
        }
    }
    & \\ \\

    \theutterance \stepcounter{utterance}  
    & & & \multicolumn{4}{p{0.6\linewidth}}{
        \cellcolor[rgb]{0.9,0.9,0.9}{
            \makecell[{{p{\linewidth}}}]{
                \texttt{\tiny{[GM$\langle$P2]}}
                \texttt{COMUNICA: Grazie per l'aggiornamento! Sulla mia griglia, il mio 'I' è stato spostato nella posizione (2,4) e il mio 'W' è stato spostato nella posizione (5,5). Sembra che i nostri oggetti 'I' e 'W' non siano ancora allineati correttamente. Vorrei discutere ulteriormente le mosse necessarie per allineare questi oggetti. Come pensi di procedere per allineare gli oggetti 'I' e 'W'?} \\
            }
        }
    }
     \\ \\

    \theutterance \stepcounter{utterance}  
    & & & \multicolumn{2}{p{0.3\linewidth}}{
        \cellcolor[rgb]{0.9,0.9,0.9}{
            \makecell[{{p{\linewidth}}}]{
                \texttt{\tiny{[GM$|$GM]}}
                \texttt{COMUNICA: Grazie per l'aggiornamento! Sulla mia griglia, il mio 'I' è stato spostato nella posizione (2,4) e il mio 'W' è stato spostato nella posizione (5,5). Sembra che i nostri oggetti 'I' e 'W' non siano ancora allineati correttamente. Vorrei discutere ulteriormente le mosse necessarie per allineare questi oggetti. Come pensi di procedere per allineare gli oggetti 'I' e 'W'?} \\
            }
        }
    }
    & & \\ \\

    \theutterance \stepcounter{utterance}  
    & & \multicolumn{4}{p{0.6\linewidth}}{
        \cellcolor[rgb]{0.9,0.9,0.9}{
            \makecell[{{p{\linewidth}}}]{
                \texttt{\tiny{[P1$\langle$GM]}}
                \texttt{Il tuo messaggio è stato inoltrato all'altro giocatore.} \\
\\ 
\\ 
\texttt{You are currently playing round 11 of maximum 20.} \\
\texttt{You have collectively accumulated 9 of 12 penalties.} \\
\texttt{L'altro giocatore ha inviato questo messaggio:} \\
\\ 
\\ 
\texttt{"Grazie per l'aggiornamento! Sulla mia griglia, il mio 'I' è stato spostato nella posizione (2,4) e il mio 'W' è stato spostato nella posizione (5,5). Sembra che i nostri oggetti 'I' e 'W' non siano ancora allineati correttamente. Vorrei discutere ulteriormente le mosse necessarie per allineare questi oggetti. Come pensi di procedere per allineare gli oggetti 'I' e 'W'?"} \\
\\ 
\\ 
\texttt{Qual è il tuo prossimo comando?} \\
            }
        }
    }
    & \\ \\

    \theutterance \stepcounter{utterance}  
    & \multicolumn{4}{p{0.6\linewidth}}{
        \cellcolor[rgb]{0.9,0.9,0.9}{
            \makecell[{{p{\linewidth}}}]{
                \texttt{\tiny{[P1$\rangle$GM]}}
                \texttt{COMUNICA: Sono d'accordo che gli oggetti 'I' e 'W' non siano ancora allineati correttamente. Considerando le posizioni attuali, potrei spostare il mio 'I' dalla posizione (2,5) alla posizione (2,4) per allinearli con il tuo 'I'. Per quanto riguarda gli oggetti 'W', potrei spostare il mio 'W' dalla posizione (4,3) alla posizione (5,5) per allinearli con il tuo 'W'. Come pensi di procedere con gli oggetti 'I' e 'W' sulla tua griglia per raggiungere l'allineamento?} \\
            }
        }
    }
    & & \\ \\

    \theutterance \stepcounter{utterance}  
    & & & \multicolumn{2}{p{0.3\linewidth}}{
        \cellcolor[rgb]{0.9,0.9,0.9}{
            \makecell[{{p{\linewidth}}}]{
                \texttt{\tiny{[GM$|$GM]}}
                \texttt{COMUNICA: Sono d'accordo che gli oggetti 'I' e 'W' non siano ancora allineati correttamente. Considerando le posizioni attuali, potrei spostare il mio 'I' dalla posizione (2,5) alla posizione (2,4) per allinearli con il tuo 'I'. Per quanto riguarda gli oggetti 'W', potrei spostare il mio 'W' dalla posizione (4,3) alla posizione (5,5) per allinearli con il tuo 'W'. Come pensi di procedere con gli oggetti 'I' e 'W' sulla tua griglia per raggiungere l'allineamento?} \\
            }
        }
    }
    & & \\ \\

    \theutterance \stepcounter{utterance}  
    & & \multicolumn{4}{p{0.6\linewidth}}{
        \cellcolor[rgb]{0.9,0.9,0.9}{
            \makecell[{{p{\linewidth}}}]{
                \texttt{\tiny{[GM$\rangle$P2]}}
                \texttt{Il tuo messaggio è stato inoltrato all'altro giocatore.} \\
\\ 
\\ 
\texttt{You are currently playing round 12 of maximum 20.} \\
\texttt{You have collectively accumulated 9 of 12 penalties.} \\
\texttt{L'altro giocatore ha inviato questo messaggio:} \\
\\ 
\\ 
\texttt{"Sono d'accordo che gli oggetti 'I' e 'W' non siano ancora allineati correttamente. Considerando le posizioni attuali, potrei spostare il mio 'I' dalla posizione (2,5) alla posizione (2,4) per allinearli con il tuo 'I'. Per quanto riguarda gli oggetti 'W', potrei spostare il mio 'W' dalla posizione (4,3) alla posizione (5,5) per allinearli con il tuo 'W'. Come pensi di procedere con gli oggetti 'I' e 'W' sulla tua griglia per raggiungere l'allineamento?"} \\
\\ 
\\ 
\texttt{Qual è il tuo prossimo comando?} \\
            }
        }
    }
    & \\ \\

    \theutterance \stepcounter{utterance}  
    & & & \multicolumn{4}{p{0.6\linewidth}}{
        \cellcolor[rgb]{0.9,0.9,0.9}{
            \makecell[{{p{\linewidth}}}]{
                \texttt{\tiny{[GM$\langle$P2]}}
                \texttt{COMUNICA: Sono d'accordo con le mosse che hai proposto per allineare gli oggetti 'I' e 'W'. Tuttavia, dal momento che il mio 'I' è già nella posizione (2,4) e il mio 'W' è già nella posizione (5,5), non è necessario che io sposti questi oggetti. Potremmo considerare di verificare se gli oggetti 'C' sono allineati correttamente o se è necessario spostarli. Come stanno gli oggetti 'C' sulla tua griglia? Sono nelle stesse posizioni degli oggetti 'C' sulla mia griglia?} \\
            }
        }
    }
     \\ \\

    \theutterance \stepcounter{utterance}  
    & & & \multicolumn{2}{p{0.3\linewidth}}{
        \cellcolor[rgb]{0.9,0.9,0.9}{
            \makecell[{{p{\linewidth}}}]{
                \texttt{\tiny{[GM$|$GM]}}
                \texttt{COMUNICA: Sono d'accordo con le mosse che hai proposto per allineare gli oggetti 'I' e 'W'. Tuttavia, dal momento che il mio 'I' è già nella posizione (2,4) e il mio 'W' è già nella posizione (5,5), non è necessario che io sposti questi oggetti. Potremmo considerare di verificare se gli oggetti 'C' sono allineati correttamente o se è necessario spostarli. Come stanno gli oggetti 'C' sulla tua griglia? Sono nelle stesse posizioni degli oggetti 'C' sulla mia griglia?} \\
            }
        }
    }
    & & \\ \\

    \theutterance \stepcounter{utterance}  
    & & \multicolumn{4}{p{0.6\linewidth}}{
        \cellcolor[rgb]{0.9,0.9,0.9}{
            \makecell[{{p{\linewidth}}}]{
                \texttt{\tiny{[P1$\langle$GM]}}
                \texttt{Il tuo messaggio è stato inoltrato all'altro giocatore.} \\
\\ 
\\ 
\texttt{You are currently playing round 12 of maximum 20.} \\
\texttt{You have collectively accumulated 9 of 12 penalties.} \\
\texttt{L'altro giocatore ha inviato questo messaggio:} \\
\\ 
\\ 
\texttt{"Sono d'accordo con le mosse che hai proposto per allineare gli oggetti 'I' e 'W'. Tuttavia, dal momento che il mio 'I' è già nella posizione (2,4) e il mio 'W' è già nella posizione (5,5), non è necessario che io sposti questi oggetti. Potremmo considerare di verificare se gli oggetti 'C' sono allineati correttamente o se è necessario spostarli. Come stanno gli oggetti 'C' sulla tua griglia? Sono nelle stesse posizioni degli oggetti 'C' sulla mia griglia?"} \\
\\ 
\\ 
\texttt{Qual è il tuo prossimo comando?} \\
            }
        }
    }
    & \\ \\

    \theutterance \stepcounter{utterance}  
    & \multicolumn{4}{p{0.6\linewidth}}{
        \cellcolor[rgb]{0.9,0.9,0.9}{
            \makecell[{{p{\linewidth}}}]{
                \texttt{\tiny{[P1$\rangle$GM]}}
                \texttt{COMUNICA: Il mio 'C' è attualmente nella posizione (2,5), ma dal momento che ho spostato il mio 'I' nella posizione (2,5) in una mossa precedente, il mio 'C' è stato spostato e ora si trova nella posizione originale. Per verificare la posizione esatta del mio 'C', devo controllare la mia griglia. Il mio 'C' è nella posizione (3,5). Come stanno gli oggetti 'C' sulla tua griglia? Sono nelle stesse posizioni del mio 'C'?} \\
            }
        }
    }
    & & \\ \\

    \theutterance \stepcounter{utterance}  
    & & & \multicolumn{2}{p{0.3\linewidth}}{
        \cellcolor[rgb]{0.9,0.9,0.9}{
            \makecell[{{p{\linewidth}}}]{
                \texttt{\tiny{[GM$|$GM]}}
                \texttt{COMUNICA: Il mio 'C' è attualmente nella posizione (2,5), ma dal momento che ho spostato il mio 'I' nella posizione (2,5) in una mossa precedente, il mio 'C' è stato spostato e ora si trova nella posizione originale. Per verificare la posizione esatta del mio 'C', devo controllare la mia griglia. Il mio 'C' è nella posizione (3,5). Come stanno gli oggetti 'C' sulla tua griglia? Sono nelle stesse posizioni del mio 'C'?} \\
            }
        }
    }
    & & \\ \\

    \theutterance \stepcounter{utterance}  
    & & \multicolumn{4}{p{0.6\linewidth}}{
        \cellcolor[rgb]{0.9,0.9,0.9}{
            \makecell[{{p{\linewidth}}}]{
                \texttt{\tiny{[GM$\rangle$P2]}}
                \texttt{Il tuo messaggio è stato inoltrato all'altro giocatore.} \\
\\ 
\\ 
\texttt{You are currently playing round 13 of maximum 20.} \\
\texttt{You have collectively accumulated 9 of 12 penalties.} \\
\texttt{L'altro giocatore ha inviato questo messaggio:} \\
\\ 
\\ 
\texttt{"Il mio 'C' è attualmente nella posizione (2,5), ma dal momento che ho spostato il mio 'I' nella posizione (2,5) in una mossa precedente, il mio 'C' è stato spostato e ora si trova nella posizione originale. Per verificare la posizione esatta del mio 'C', devo controllare la mia griglia. Il mio 'C' è nella posizione (3,5). Come stanno gli oggetti 'C' sulla tua griglia? Sono nelle stesse posizioni del mio 'C'?"} \\
\\ 
\\ 
\texttt{Qual è il tuo prossimo comando?} \\
            }
        }
    }
    & \\ \\

    \theutterance \stepcounter{utterance}  
    & & & \multicolumn{4}{p{0.6\linewidth}}{
        \cellcolor[rgb]{0.9,0.9,0.9}{
            \makecell[{{p{\linewidth}}}]{
                \texttt{\tiny{[GM$\langle$P2]}}
                \texttt{COMUNICA: Il mio 'C' è nella posizione (4,2). Sembra che gli oggetti 'C' non siano allineati correttamente. Considerando le posizioni attuali, potremmo discutere le mosse necessarie per allineare gli oggetti 'C'. Tuttavia, prima di procedere, vorrei assicurarmi che tutti gli altri oggetti siano nella posizione corretta. Come stanno gli oggetti 'H', 'I', 'T' e 'W' sulla tua griglia? Sono tutte nella posizione che abbiamo discusso e pianificato?} \\
            }
        }
    }
     \\ \\

    \theutterance \stepcounter{utterance}  
    & & & \multicolumn{2}{p{0.3\linewidth}}{
        \cellcolor[rgb]{0.9,0.9,0.9}{
            \makecell[{{p{\linewidth}}}]{
                \texttt{\tiny{[GM$|$GM]}}
                \texttt{COMUNICA: Il mio 'C' è nella posizione (4,2). Sembra che gli oggetti 'C' non siano allineati correttamente. Considerando le posizioni attuali, potremmo discutere le mosse necessarie per allineare gli oggetti 'C'. Tuttavia, prima di procedere, vorrei assicurarmi che tutti gli altri oggetti siano nella posizione corretta. Come stanno gli oggetti 'H', 'I', 'T' e 'W' sulla tua griglia? Sono tutte nella posizione che abbiamo discusso e pianificato?} \\
            }
        }
    }
    & & \\ \\

    \theutterance \stepcounter{utterance}  
    & & \multicolumn{4}{p{0.6\linewidth}}{
        \cellcolor[rgb]{0.9,0.9,0.9}{
            \makecell[{{p{\linewidth}}}]{
                \texttt{\tiny{[P1$\langle$GM]}}
                \texttt{Il tuo messaggio è stato inoltrato all'altro giocatore.} \\
\\ 
\\ 
\texttt{You are currently playing round 13 of maximum 20.} \\
\texttt{You have collectively accumulated 9 of 12 penalties.} \\
\texttt{L'altro giocatore ha inviato questo messaggio:} \\
\\ 
\\ 
\texttt{"Il mio 'C' è nella posizione (4,2). Sembra che gli oggetti 'C' non siano allineati correttamente. Considerando le posizioni attuali, potremmo discutere le mosse necessarie per allineare gli oggetti 'C'. Tuttavia, prima di procedere, vorrei assicurarmi che tutti gli altri oggetti siano nella posizione corretta. Come stanno gli oggetti 'H', 'I', 'T' e 'W' sulla tua griglia? Sono tutte nella posizione che abbiamo discusso e pianificato?"} \\
\\ 
\\ 
\texttt{Qual è il tuo prossimo comando?} \\
            }
        }
    }
    & \\ \\

    \theutterance \stepcounter{utterance}  
    & \multicolumn{4}{p{0.6\linewidth}}{
        \cellcolor[rgb]{0.9,0.9,0.9}{
            \makecell[{{p{\linewidth}}}]{
                \texttt{\tiny{[P1$\rangle$GM]}}
                \texttt{COMUNICA: Gli oggetti 'H' è nella posizione (3,1), l'oggetto 'I' è nella posizione (2,5), l'oggetto 'T' è nella posizione (3,7) e l'oggetto 'W' è nella posizione (4,3). Sembra che gli oggetti 'I' e 'W' non siano ancora stati spostati nelle posizioni finali che abbiamo discusso. Ricordiamoci di spostare il mio 'I' nella posizione (2,4) e il mio 'W' nella posizione (5,5) per allinearli con i tuoi oggetti 'I' e 'W'. Dopo aver verificato, possiamo procedere con gli oggetti 'C' per assicurarci che siano allineati correttamente. Come procediamo con gli oggetti 'C' e gli altri oggetti che non sono ancora nella posizione corretta?} \\
            }
        }
    }
    & & \\ \\

    \theutterance \stepcounter{utterance}  
    & & & \multicolumn{2}{p{0.3\linewidth}}{
        \cellcolor[rgb]{0.9,0.9,0.9}{
            \makecell[{{p{\linewidth}}}]{
                \texttt{\tiny{[GM$|$GM]}}
                \texttt{COMUNICA: Gli oggetti 'H' è nella posizione (3,1), l'oggetto 'I' è nella posizione (2,5), l'oggetto 'T' è nella posizione (3,7) e l'oggetto 'W' è nella posizione (4,3). Sembra che gli oggetti 'I' e 'W' non siano ancora stati spostati nelle posizioni finali che abbiamo discusso. Ricordiamoci di spostare il mio 'I' nella posizione (2,4) e il mio 'W' nella posizione (5,5) per allinearli con i tuoi oggetti 'I' e 'W'. Dopo aver verificato, possiamo procedere con gli oggetti 'C' per assicurarci che siano allineati correttamente. Come procediamo con gli oggetti 'C' e gli altri oggetti che non sono ancora nella posizione corretta?} \\
            }
        }
    }
    & & \\ \\

    \theutterance \stepcounter{utterance}  
    & & \multicolumn{4}{p{0.6\linewidth}}{
        \cellcolor[rgb]{0.9,0.9,0.9}{
            \makecell[{{p{\linewidth}}}]{
                \texttt{\tiny{[GM$\rangle$P2]}}
                \texttt{Il tuo messaggio è stato inoltrato all'altro giocatore.} \\
\\ 
\\ 
\texttt{You are currently playing round 14 of maximum 20.} \\
\texttt{You have collectively accumulated 9 of 12 penalties.} \\
\texttt{L'altro giocatore ha inviato questo messaggio:} \\
\\ 
\\ 
\texttt{"Gli oggetti 'H' è nella posizione (3,1), l'oggetto 'I' è nella posizione (2,5), l'oggetto 'T' è nella posizione (3,7) e l'oggetto 'W' è nella posizione (4,3). Sembra che gli oggetti 'I' e 'W' non siano ancora stati spostati nelle posizioni finali che abbiamo discusso. Ricordiamoci di spostare il mio 'I' nella posizione (2,4) e il mio 'W' nella posizione (5,5) per allinearli con i tuoi oggetti 'I' e 'W'. Dopo aver verificato, possiamo procedere con gli oggetti 'C' per assicurarci che siano allineati correttamente. Come procediamo con gli oggetti 'C' e gli altri oggetti che non sono ancora nella posizione corretta?"} \\
\\ 
\\ 
\texttt{Qual è il tuo prossimo comando?} \\
            }
        }
    }
    & \\ \\

    \theutterance \stepcounter{utterance}  
    & & & \multicolumn{4}{p{0.6\linewidth}}{
        \cellcolor[rgb]{0.9,0.9,0.9}{
            \makecell[{{p{\linewidth}}}]{
                \texttt{\tiny{[GM$\langle$P2]}}
                \texttt{COMUNICA: Sono d'accordo che gli oggetti 'I' e 'W' non siano ancora stati spostati nelle posizioni finali. Tuttavia, ricordo che il mio 'I' è già nella posizione (2,4) e il mio 'W' è già nella posizione (5,5). Quindi, non è necessario che io sposti questi oggetti. Per quanto riguarda gli oggetti 'C', potremmo discutere le mosse necessarie per allineare gli oggetti 'C'. Poiché il mio 'C' è nella posizione (4,2) e il tuo 'C' è nella posizione (3,5), potremmo trovare una posizione comune per gli oggetti 'C'. Come pensi di procedere con gli oggetti 'C'?} \\
            }
        }
    }
     \\ \\

    \theutterance \stepcounter{utterance}  
    & & & \multicolumn{2}{p{0.3\linewidth}}{
        \cellcolor[rgb]{0.9,0.9,0.9}{
            \makecell[{{p{\linewidth}}}]{
                \texttt{\tiny{[GM$|$GM]}}
                \texttt{COMUNICA: Sono d'accordo che gli oggetti 'I' e 'W' non siano ancora stati spostati nelle posizioni finali. Tuttavia, ricordo che il mio 'I' è già nella posizione (2,4) e il mio 'W' è già nella posizione (5,5). Quindi, non è necessario che io sposti questi oggetti. Per quanto riguarda gli oggetti 'C', potremmo discutere le mosse necessarie per allineare gli oggetti 'C'. Poiché il mio 'C' è nella posizione (4,2) e il tuo 'C' è nella posizione (3,5), potremmo trovare una posizione comune per gli oggetti 'C'. Come pensi di procedere con gli oggetti 'C'?} \\
            }
        }
    }
    & & \\ \\

    \theutterance \stepcounter{utterance}  
    & & \multicolumn{4}{p{0.6\linewidth}}{
        \cellcolor[rgb]{0.9,0.9,0.9}{
            \makecell[{{p{\linewidth}}}]{
                \texttt{\tiny{[P1$\langle$GM]}}
                \texttt{Il tuo messaggio è stato inoltrato all'altro giocatore.} \\
\\ 
\\ 
\texttt{You are currently playing round 14 of maximum 20.} \\
\texttt{You have collectively accumulated 9 of 12 penalties.} \\
\texttt{L'altro giocatore ha inviato questo messaggio:} \\
\\ 
\\ 
\texttt{"Sono d'accordo che gli oggetti 'I' e 'W' non siano ancora stati spostati nelle posizioni finali. Tuttavia, ricordo che il mio 'I' è già nella posizione (2,4) e il mio 'W' è già nella posizione (5,5). Quindi, non è necessario che io sposti questi oggetti. Per quanto riguarda gli oggetti 'C', potremmo discutere le mosse necessarie per allineare gli oggetti 'C'. Poiché il mio 'C' è nella posizione (4,2) e il tuo 'C' è nella posizione (3,5), potremmo trovare una posizione comune per gli oggetti 'C'. Come pensi di procedere con gli oggetti 'C'?"} \\
\\ 
\\ 
\texttt{Qual è il tuo prossimo comando?} \\
            }
        }
    }
    & \\ \\

    \theutterance \stepcounter{utterance}  
    & \multicolumn{4}{p{0.6\linewidth}}{
        \cellcolor[rgb]{0.9,0.9,0.9}{
            \makecell[{{p{\linewidth}}}]{
                \texttt{\tiny{[P1$\rangle$GM]}}
                \texttt{COMUNICA: Sono d'accordo che il tuo 'I' e 'W' sono già nelle posizioni finali. Per gli oggetti 'C', considerando le posizioni attuali, potremmo scegliere una posizione intermedia per allineare gli oggetti 'C'. Una possibile opzione potrebbe essere la posizione (4,5) o una posizione vicina che sia accessibile per entrambi. Tuttavia, dobbiamo anche considerare le posizioni degli altri oggetti e assicurare che non ci siano ostacoli per spostare gli oggetti 'C' nella nuova posizione. Come pensi di procedere con gli oggetti 'C' e quali posizioni Alternate potremmo considerare?} \\
            }
        }
    }
    & & \\ \\

    \theutterance \stepcounter{utterance}  
    & & & \multicolumn{2}{p{0.3\linewidth}}{
        \cellcolor[rgb]{0.9,0.9,0.9}{
            \makecell[{{p{\linewidth}}}]{
                \texttt{\tiny{[GM$|$GM]}}
                \texttt{COMUNICA: Sono d'accordo che il tuo 'I' e 'W' sono già nelle posizioni finali. Per gli oggetti 'C', considerando le posizioni attuali, potremmo scegliere una posizione intermedia per allineare gli oggetti 'C'. Una possibile opzione potrebbe essere la posizione (4,5) o una posizione vicina che sia accessibile per entrambi. Tuttavia, dobbiamo anche considerare le posizioni degli altri oggetti e assicurare che non ci siano ostacoli per spostare gli oggetti 'C' nella nuova posizione. Come pensi di procedere con gli oggetti 'C' e quali posizioni Alternate potremmo considerare?} \\
            }
        }
    }
    & & \\ \\

    \theutterance \stepcounter{utterance}  
    & & \multicolumn{4}{p{0.6\linewidth}}{
        \cellcolor[rgb]{0.9,0.9,0.9}{
            \makecell[{{p{\linewidth}}}]{
                \texttt{\tiny{[GM$\rangle$P2]}}
                \texttt{Il tuo messaggio è stato inoltrato all'altro giocatore.} \\
\\ 
\\ 
\texttt{You are currently playing round 15 of maximum 20.} \\
\texttt{You have collectively accumulated 9 of 12 penalties.} \\
\texttt{L'altro giocatore ha inviato questo messaggio:} \\
\\ 
\\ 
\texttt{"Sono d'accordo che il tuo 'I' e 'W' sono già nelle posizioni finali. Per gli oggetti 'C', considerando le posizioni attuali, potremmo scegliere una posizione intermedia per allineare gli oggetti 'C'. Una possibile opzione potrebbe essere la posizione (4,5) o una posizione vicina che sia accessibile per entrambi. Tuttavia, dobbiamo anche considerare le posizioni degli altri oggetti e assicurare che non ci siano ostacoli per spostare gli oggetti 'C' nella nuova posizione. Come pensi di procedere con gli oggetti 'C' e quali posizioni Alternate potremmo considerare?"} \\
\\ 
\\ 
\texttt{Qual è il tuo prossimo comando?} \\
            }
        }
    }
    & \\ \\

    \theutterance \stepcounter{utterance}  
    & & & \multicolumn{4}{p{0.6\linewidth}}{
        \cellcolor[rgb]{0.9,0.9,0.9}{
            \makecell[{{p{\linewidth}}}]{
                \texttt{\tiny{[GM$\langle$P2]}}
                \texttt{COMUNICA: Sono d'accordo che una posizione intermedia per gli oggetti 'C' potrebbe essere una buona opzione. La posizione (4,5) sembra ragionevole, ma dobbiamo considerare gli ostacoli sulla griglia. Notando che la mia griglia ha una struttura con linee e angoli, potremmo dover trovare un percorso alternativo per spostare gli oggetti 'C' nella nuova posizione. Una possibile alternativa potrebbe essere la posizione (3,2) o (3,3), che potrebbe essere più accessibile per entrambi. Come pensi di procedere con gli oggetti 'C' e quali percorsi possiamo considerare per spostarli nella nuova posizione?} \\
            }
        }
    }
     \\ \\

    \theutterance \stepcounter{utterance}  
    & & & \multicolumn{2}{p{0.3\linewidth}}{
        \cellcolor[rgb]{0.9,0.9,0.9}{
            \makecell[{{p{\linewidth}}}]{
                \texttt{\tiny{[GM$|$GM]}}
                \texttt{COMUNICA: Sono d'accordo che una posizione intermedia per gli oggetti 'C' potrebbe essere una buona opzione. La posizione (4,5) sembra ragionevole, ma dobbiamo considerare gli ostacoli sulla griglia. Notando che la mia griglia ha una struttura con linee e angoli, potremmo dover trovare un percorso alternativo per spostare gli oggetti 'C' nella nuova posizione. Una possibile alternativa potrebbe essere la posizione (3,2) o (3,3), che potrebbe essere più accessibile per entrambi. Come pensi di procedere con gli oggetti 'C' e quali percorsi possiamo considerare per spostarli nella nuova posizione?} \\
            }
        }
    }
    & & \\ \\

    \theutterance \stepcounter{utterance}  
    & & \multicolumn{4}{p{0.6\linewidth}}{
        \cellcolor[rgb]{0.9,0.9,0.9}{
            \makecell[{{p{\linewidth}}}]{
                \texttt{\tiny{[P1$\langle$GM]}}
                \texttt{Il tuo messaggio è stato inoltrato all'altro giocatore.} \\
\\ 
\\ 
\texttt{You are currently playing round 15 of maximum 20.} \\
\texttt{You have collectively accumulated 9 of 12 penalties.} \\
\texttt{L'altro giocatore ha inviato questo messaggio:} \\
\\ 
\\ 
\texttt{"Sono d'accordo che una posizione intermedia per gli oggetti 'C' potrebbe essere una buona opzione. La posizione (4,5) sembra ragionevole, ma dobbiamo considerare gli ostacoli sulla griglia. Notando che la mia griglia ha una struttura con linee e angoli, potremmo dover trovare un percorso alternativo per spostare gli oggetti 'C' nella nuova posizione. Una possibile alternativa potrebbe essere la posizione (3,2) o (3,3), che potrebbe essere più accessibile per entrambi. Come pensi di procedere con gli oggetti 'C' e quali percorsi possiamo considerare per spostarli nella nuova posizione?"} \\
\\ 
\\ 
\texttt{Qual è il tuo prossimo comando?} \\
            }
        }
    }
    & \\ \\

    \theutterance \stepcounter{utterance}  
    & \multicolumn{4}{p{0.6\linewidth}}{
        \cellcolor[rgb]{0.9,0.9,0.9}{
            \makecell[{{p{\linewidth}}}]{
                \texttt{\tiny{[P1$\rangle$GM]}}
                \texttt{COMUNICA: Sono d'accordo che la posizione (4,5) potrebbe essere difficile da raggiungere a causa degli ostacoli sulla griglia. Le posizioni alternative che hai proposto, (3,2) o (3,3), sembrano più accessibili. Tuttavia, dobbiamo anche considerare la posizione del mio 'C' attuale, che è (3,5), e la posizione del tuo 'C', che è (4,2). Forse potremmo trovare una posizione che sia più vicina a entrambi, ad esempio (3,4) o (4,3), se sono accessibili. Quali sono le tue thought su queste posizioni e come possiamo procedere per spostare gli oggetti 'C' nella nuova posizione?} \\
            }
        }
    }
    & & \\ \\

    \theutterance \stepcounter{utterance}  
    & & & \multicolumn{2}{p{0.3\linewidth}}{
        \cellcolor[rgb]{0.9,0.9,0.9}{
            \makecell[{{p{\linewidth}}}]{
                \texttt{\tiny{[GM$|$GM]}}
                \texttt{COMUNICA: Sono d'accordo che la posizione (4,5) potrebbe essere difficile da raggiungere a causa degli ostacoli sulla griglia. Le posizioni alternative che hai proposto, (3,2) o (3,3), sembrano più accessibili. Tuttavia, dobbiamo anche considerare la posizione del mio 'C' attuale, che è (3,5), e la posizione del tuo 'C', che è (4,2). Forse potremmo trovare una posizione che sia più vicina a entrambi, ad esempio (3,4) o (4,3), se sono accessibili. Quali sono le tue thought su queste posizioni e come possiamo procedere per spostare gli oggetti 'C' nella nuova posizione?} \\
            }
        }
    }
    & & \\ \\

    \theutterance \stepcounter{utterance}  
    & & \multicolumn{4}{p{0.6\linewidth}}{
        \cellcolor[rgb]{0.9,0.9,0.9}{
            \makecell[{{p{\linewidth}}}]{
                \texttt{\tiny{[GM$\rangle$P2]}}
                \texttt{Il tuo messaggio è stato inoltrato all'altro giocatore.} \\
\\ 
\\ 
\texttt{You are currently playing round 16 of maximum 20.} \\
\texttt{You have collectively accumulated 9 of 12 penalties.} \\
\texttt{L'altro giocatore ha inviato questo messaggio:} \\
\\ 
\\ 
\texttt{"Sono d'accordo che la posizione (4,5) potrebbe essere difficile da raggiungere a causa degli ostacoli sulla griglia. Le posizioni alternative che hai proposto, (3,2) o (3,3), sembrano più accessibili. Tuttavia, dobbiamo anche considerare la posizione del mio 'C' attuale, che è (3,5), e la posizione del tuo 'C', che è (4,2). Forse potremmo trovare una posizione che sia più vicina a entrambi, ad esempio (3,4) o (4,3), se sono accessibili. Quali sono le tue thought su queste posizioni e come possiamo procedere per spostare gli oggetti 'C' nella nuova posizione?"} \\
\\ 
\\ 
\texttt{Qual è il tuo prossimo comando?} \\
            }
        }
    }
    & \\ \\

    \theutterance \stepcounter{utterance}  
    & & & \multicolumn{4}{p{0.6\linewidth}}{
        \cellcolor[rgb]{0.9,0.9,0.9}{
            \makecell[{{p{\linewidth}}}]{
                \texttt{\tiny{[GM$\langle$P2]}}
                \texttt{COMUNICA: Sono d'accordo che le posizioni (3,4) o (4,3) potrebbero essere più vicine a entrambi gli oggetti 'C' e quindi più accessibili. Tuttavia, dobbiamo anche considerare gli ostacoli sulla griglia e gli altri oggetti che potrebbero bloccare il percorso. La posizione (3,4) sembra una buona opzione, ma dobbiamo verificare se è libera e se possiamo spostare gli oggetti 'C' in quella posizione senza incontrare ostacoli. Una volta che avremo confermato la possibilità di spostare gli oggetti 'C' in quella posizione, potremmo procedere con la mossa. Come pensi di procedere per verificare la posizione (3,4) e spostare gli oggetti 'C'?} \\
            }
        }
    }
     \\ \\

    \theutterance \stepcounter{utterance}  
    & & & \multicolumn{2}{p{0.3\linewidth}}{
        \cellcolor[rgb]{0.9,0.9,0.9}{
            \makecell[{{p{\linewidth}}}]{
                \texttt{\tiny{[GM$|$GM]}}
                \texttt{COMUNICA: Sono d'accordo che le posizioni (3,4) o (4,3) potrebbero essere più vicine a entrambi gli oggetti 'C' e quindi più accessibili. Tuttavia, dobbiamo anche considerare gli ostacoli sulla griglia e gli altri oggetti che potrebbero bloccare il percorso. La posizione (3,4) sembra una buona opzione, ma dobbiamo verificare se è libera e se possiamo spostare gli oggetti 'C' in quella posizione senza incontrare ostacoli. Una volta che avremo confermato la possibilità di spostare gli oggetti 'C' in quella posizione, potremmo procedere con la mossa. Come pensi di procedere per verificare la posizione (3,4) e spostare gli oggetti 'C'?} \\
            }
        }
    }
    & & \\ \\

    \theutterance \stepcounter{utterance}  
    & & \multicolumn{4}{p{0.6\linewidth}}{
        \cellcolor[rgb]{0.9,0.9,0.9}{
            \makecell[{{p{\linewidth}}}]{
                \texttt{\tiny{[P1$\langle$GM]}}
                \texttt{Il tuo messaggio è stato inoltrato all'altro giocatore.} \\
\\ 
\\ 
\texttt{You are currently playing round 16 of maximum 20.} \\
\texttt{You have collectively accumulated 9 of 12 penalties.} \\
\texttt{L'altro giocatore ha inviato questo messaggio:} \\
\\ 
\\ 
\texttt{"Sono d'accordo che le posizioni (3,4) o (4,3) potrebbero essere più vicine a entrambi gli oggetti 'C' e quindi più accessibili. Tuttavia, dobbiamo anche considerare gli ostacoli sulla griglia e gli altri oggetti che potrebbero bloccare il percorso. La posizione (3,4) sembra una buona opzione, ma dobbiamo verificare se è libera e se possiamo spostare gli oggetti 'C' in quella posizione senza incontrare ostacoli. Una volta che avremo confermato la possibilità di spostare gli oggetti 'C' in quella posizione, potremmo procedere con la mossa. Come pensi di procedere per verificare la posizione (3,4) e spostare gli oggetti 'C'?"} \\
\\ 
\\ 
\texttt{Qual è il tuo prossimo comando?} \\
            }
        }
    }
    & \\ \\

    \theutterance \stepcounter{utterance}  
    & \multicolumn{4}{p{0.6\linewidth}}{
        \cellcolor[rgb]{0.9,0.9,0.9}{
            \makecell[{{p{\linewidth}}}]{
                \texttt{\tiny{[P1$\rangle$GM]}}
                \texttt{COMUNICA: Sono d'accordo che la posizione (3,4) sembra una buona opzione, ma dobbiamo verificare se è libera e accessibile. Per verificare la posizione, potrei esaminare la mia griglia e verificare se ci sono ostacoli o altri oggetti che potrebbero bloccare il percorso. Se la posizione (3,4) è libera, potrei spostare il mio 'C' in quella posizione. Intanto, potresti verificare la tua griglia e vedere se puoi spostare il tuo 'C' nella stessa posizione. Se entrambi confermiamo che la posizione (3,4) è accessibile, potremmo procedere con la mossa e spostare gli oggetti 'C' in quella posizione. Come pensi di procedere per verificare la posizione (3,4) sulla tua griglia?} \\
            }
        }
    }
    & & \\ \\

    \theutterance \stepcounter{utterance}  
    & & & \multicolumn{2}{p{0.3\linewidth}}{
        \cellcolor[rgb]{0.9,0.9,0.9}{
            \makecell[{{p{\linewidth}}}]{
                \texttt{\tiny{[GM$|$GM]}}
                \texttt{COMUNICA: Sono d'accordo che la posizione (3,4) sembra una buona opzione, ma dobbiamo verificare se è libera e accessibile. Per verificare la posizione, potrei esaminare la mia griglia e verificare se ci sono ostacoli o altri oggetti che potrebbero bloccare il percorso. Se la posizione (3,4) è libera, potrei spostare il mio 'C' in quella posizione. Intanto, potresti verificare la tua griglia e vedere se puoi spostare il tuo 'C' nella stessa posizione. Se entrambi confermiamo che la posizione (3,4) è accessibile, potremmo procedere con la mossa e spostare gli oggetti 'C' in quella posizione. Come pensi di procedere per verificare la posizione (3,4) sulla tua griglia?} \\
            }
        }
    }
    & & \\ \\

    \theutterance \stepcounter{utterance}  
    & & \multicolumn{4}{p{0.6\linewidth}}{
        \cellcolor[rgb]{0.9,0.9,0.9}{
            \makecell[{{p{\linewidth}}}]{
                \texttt{\tiny{[GM$\rangle$P2]}}
                \texttt{Il tuo messaggio è stato inoltrato all'altro giocatore.} \\
\\ 
\\ 
\texttt{You are currently playing round 17 of maximum 20.} \\
\texttt{You have collectively accumulated 9 of 12 penalties.} \\
\texttt{L'altro giocatore ha inviato questo messaggio:} \\
\\ 
\\ 
\texttt{"Sono d'accordo che la posizione (3,4) sembra una buona opzione, ma dobbiamo verificare se è libera e accessibile. Per verificare la posizione, potrei esaminare la mia griglia e verificare se ci sono ostacoli o altri oggetti che potrebbero bloccare il percorso. Se la posizione (3,4) è libera, potrei spostare il mio 'C' in quella posizione. Intanto, potresti verificare la tua griglia e vedere se puoi spostare il tuo 'C' nella stessa posizione. Se entrambi confermiamo che la posizione (3,4) è accessibile, potremmo procedere con la mossa e spostare gli oggetti 'C' in quella posizione. Come pensi di procedere per verificare la posizione (3,4) sulla tua griglia?"} \\
\\ 
\\ 
\texttt{Qual è il tuo prossimo comando?} \\
            }
        }
    }
    & \\ \\

    \theutterance \stepcounter{utterance}  
    & & & \multicolumn{4}{p{0.6\linewidth}}{
        \cellcolor[rgb]{0.9,0.9,0.9}{
            \makecell[{{p{\linewidth}}}]{
                \texttt{\tiny{[GM$\langle$P2]}}
                \texttt{COMUNICA: Sono d'accordo con il piano di verificare la posizione (3,4) sulla nostra griglia. Esaminerò la mia griglia per verificare se la posizione (3,4) è libera e accessibile. Se la posizione è libera, potrei spostare il mio 'C' dalla posizione (4,2) alla posizione (3,4). Spero che anche tu possa verificare la tua griglia e confermare se la posizione (3,4) è accessibile per il tuo 'C'. Se entrambi confermiamo che la posizione (3,4) è libera, potremmo procedere con la mossa e allineare gli oggetti 'C'. Procederò con la verifica e ti aggiornerò sui risultati.} \\
            }
        }
    }
     \\ \\

    \theutterance \stepcounter{utterance}  
    & & & \multicolumn{2}{p{0.3\linewidth}}{
        \cellcolor[rgb]{0.9,0.9,0.9}{
            \makecell[{{p{\linewidth}}}]{
                \texttt{\tiny{[GM$|$GM]}}
                \texttt{COMUNICA: Sono d'accordo con il piano di verificare la posizione (3,4) sulla nostra griglia. Esaminerò la mia griglia per verificare se la posizione (3,4) è libera e accessibile. Se la posizione è libera, potrei spostare il mio 'C' dalla posizione (4,2) alla posizione (3,4). Spero che anche tu possa verificare la tua griglia e confermare se la posizione (3,4) è accessibile per il tuo 'C'. Se entrambi confermiamo che la posizione (3,4) è libera, potremmo procedere con la mossa e allineare gli oggetti 'C'. Procederò con la verifica e ti aggiornerò sui risultati.} \\
            }
        }
    }
    & & \\ \\

    \theutterance \stepcounter{utterance}  
    & & \multicolumn{4}{p{0.6\linewidth}}{
        \cellcolor[rgb]{0.9,0.9,0.9}{
            \makecell[{{p{\linewidth}}}]{
                \texttt{\tiny{[P1$\langle$GM]}}
                \texttt{Il tuo messaggio è stato inoltrato all'altro giocatore.} \\
\\ 
\\ 
\texttt{You are currently playing round 17 of maximum 20.} \\
\texttt{You have collectively accumulated 9 of 12 penalties.} \\
\texttt{L'altro giocatore ha inviato questo messaggio:} \\
\\ 
\\ 
\texttt{"Sono d'accordo con il piano di verificare la posizione (3,4) sulla nostra griglia. Esaminerò la mia griglia per verificare se la posizione (3,4) è libera e accessibile. Se la posizione è libera, potrei spostare il mio 'C' dalla posizione (4,2) alla posizione (3,4). Spero che anche tu possa verificare la tua griglia e confermare se la posizione (3,4) è accessibile per il tuo 'C'. Se entrambi confermiamo che la posizione (3,4) è libera, potremmo procedere con la mossa e allineare gli oggetti 'C'. Procederò con la verifica e ti aggiornerò sui risultati."} \\
\\ 
\\ 
\texttt{Qual è il tuo prossimo comando?} \\
            }
        }
    }
    & \\ \\

    \theutterance \stepcounter{utterance}  
    & \multicolumn{4}{p{0.6\linewidth}}{
        \cellcolor[rgb]{0.9,0.9,0.9}{
            \makecell[{{p{\linewidth}}}]{
                \texttt{\tiny{[P1$\rangle$GM]}}
                \texttt{COMUNICA: Ho verificato la mia griglia e posso confermare che la posizione (3,4) è accessibile per il mio 'C'. Il mio 'C' attualmente si trova nella posizione (3,5), quindi potrei spostarlo nella posizione (3,4) senza incontrare ostacoli. Sono pronto a procedere con la mossa e allineare gli oggetti 'C' nella posizione (3,4). Aspetto la tua conferma per procedere con la mossa e completare l'allineamento degli oggetti 'C'.} \\
            }
        }
    }
    & & \\ \\

    \theutterance \stepcounter{utterance}  
    & & & \multicolumn{2}{p{0.3\linewidth}}{
        \cellcolor[rgb]{0.9,0.9,0.9}{
            \makecell[{{p{\linewidth}}}]{
                \texttt{\tiny{[GM$|$GM]}}
                \texttt{COMUNICA: Ho verificato la mia griglia e posso confermare che la posizione (3,4) è accessibile per il mio 'C'. Il mio 'C' attualmente si trova nella posizione (3,5), quindi potrei spostarlo nella posizione (3,4) senza incontrare ostacoli. Sono pronto a procedere con la mossa e allineare gli oggetti 'C' nella posizione (3,4). Aspetto la tua conferma per procedere con la mossa e completare l'allineamento degli oggetti 'C'.} \\
            }
        }
    }
    & & \\ \\

    \theutterance \stepcounter{utterance}  
    & & \multicolumn{4}{p{0.6\linewidth}}{
        \cellcolor[rgb]{0.9,0.9,0.9}{
            \makecell[{{p{\linewidth}}}]{
                \texttt{\tiny{[GM$\rangle$P2]}}
                \texttt{Il tuo messaggio è stato inoltrato all'altro giocatore.} \\
\\ 
\\ 
\texttt{You are currently playing round 18 of maximum 20.} \\
\texttt{You have collectively accumulated 9 of 12 penalties.} \\
\texttt{L'altro giocatore ha inviato questo messaggio:} \\
\\ 
\\ 
\texttt{"Ho verificato la mia griglia e posso confermare che la posizione (3,4) è accessibile per il mio 'C'. Il mio 'C' attualmente si trova nella posizione (3,5), quindi potrei spostarlo nella posizione (3,4) senza incontrare ostacoli. Sono pronto a procedere con la mossa e allineare gli oggetti 'C' nella posizione (3,4). Aspetto la tua conferma per procedere con la mossa e completare l'allineamento degli oggetti 'C'."} \\
\\ 
\\ 
\texttt{Qual è il tuo prossimo comando?} \\
            }
        }
    }
    & \\ \\

    \theutterance \stepcounter{utterance}  
    & & & \multicolumn{4}{p{0.6\linewidth}}{
        \cellcolor[rgb]{0.9,0.9,0.9}{
            \makecell[{{p{\linewidth}}}]{
                \texttt{\tiny{[GM$\langle$P2]}}
                \texttt{SPOSTA: C, (3,4)} \\
            }
        }
    }
     \\ \\

    \theutterance \stepcounter{utterance}  
    & & & \multicolumn{2}{p{0.3\linewidth}}{
        \cellcolor[rgb]{0.9,0.9,0.9}{
            \makecell[{{p{\linewidth}}}]{
                \texttt{\tiny{[GM$|$GM]}}
                \texttt{SPOSTA: C, (3,4)} \\
            }
        }
    }
    & & \\ \\

    \theutterance \stepcounter{utterance}  
    & & & \multicolumn{2}{p{0.3\linewidth}}{
        \cellcolor[rgb]{0.9,0.9,0.9}{
            \makecell[{{p{\linewidth}}}]{
                \texttt{\tiny{[GM$|$GM]}}
                \texttt{Spostato 'C' in (3,4) con successo. La tua griglia aggiornata appare così:} \\
\texttt{```} \\
\texttt{ 1234567} \\
\texttt{╔═══════╗} \\
\texttt{║◌◌H◌◌◌◌║ 1} \\
\texttt{║◌◌◌◌◌◌◌║ 2} \\
\texttt{║◌◌◌◌◌┌─╢ 3} \\
\texttt{║◌IC┌─┤◌║ 4} \\
\texttt{║◌◌◌│W├─╢ 5} \\
\texttt{╟───┼─┘◌║ 6} \\
\texttt{║◌◌T│◌◌◌║ 7} \\
\texttt{╚═══╧═══╝} \\
\texttt{```} \\
            }
        }
    }
    & & \\ \\

    \theutterance \stepcounter{utterance}  
    & & \multicolumn{4}{p{0.6\linewidth}}{
        \cellcolor[rgb]{0.9,0.9,0.9}{
            \makecell[{{p{\linewidth}}}]{
                \texttt{\tiny{[P1$\langle$GM]}}
                \texttt{Il tuo messaggio è stato inoltrato all'altro giocatore.} \\
\\ 
\\ 
\texttt{You are currently playing round 18 of maximum 20.} \\
\texttt{You have collectively accumulated 9 of 12 penalties.} \\
\texttt{L'altro giocatore ha spostato un oggetto sulla sua griglia.} \\
\\ 
\\ 
\texttt{Qual è il tuo prossimo comando?} \\
            }
        }
    }
    & \\ \\

    \theutterance \stepcounter{utterance}  
    & \multicolumn{4}{p{0.6\linewidth}}{
        \cellcolor[rgb]{0.9,0.9,0.9}{
            \makecell[{{p{\linewidth}}}]{
                \texttt{\tiny{[P1$\rangle$GM]}}
                \texttt{SPOSTA: C, (3,4)} \\
            }
        }
    }
    & & \\ \\

    \theutterance \stepcounter{utterance}  
    & & & \multicolumn{2}{p{0.3\linewidth}}{
        \cellcolor[rgb]{0.9,0.9,0.9}{
            \makecell[{{p{\linewidth}}}]{
                \texttt{\tiny{[GM$|$GM]}}
                \texttt{SPOSTA: C, (3,4)} \\
            }
        }
    }
    & & \\ \\

    \theutterance \stepcounter{utterance}  
    & & & \multicolumn{2}{p{0.3\linewidth}}{
        \cellcolor[rgb]{0.9,0.9,0.9}{
            \makecell[{{p{\linewidth}}}]{
                \texttt{\tiny{[GM$|$GM]}}
                \texttt{Spostato 'C' in (3,4) con successo. La tua griglia aggiornata appare così:} \\
\texttt{```} \\
\texttt{ 1234567} \\
\texttt{╔═══════╗} \\
\texttt{║◌◌H◌◌◌◌║ 1} \\
\texttt{║◌◌◌◌◌◌◌║ 2} \\
\texttt{║◌◌◌W◌┌─╢ 3} \\
\texttt{║◌◌C┌─┤◌║ 4} \\
\texttt{║◌I◌│◌├─╢ 5} \\
\texttt{╟───┼─┘◌║ 6} \\
\texttt{║◌◌T│◌◌◌║ 7} \\
\texttt{╚═══╧═══╝} \\
\texttt{```} \\
            }
        }
    }
    & & \\ \\

    \theutterance \stepcounter{utterance}  
    & & \multicolumn{4}{p{0.6\linewidth}}{
        \cellcolor[rgb]{0.9,0.9,0.9}{
            \makecell[{{p{\linewidth}}}]{
                \texttt{\tiny{[GM$\rangle$P2]}}
                \texttt{Spostato 'C' in (3,4) con successo. La tua griglia aggiornata appare così:} \\
\texttt{```} \\
\texttt{ 1234567} \\
\texttt{╔═══════╗} \\
\texttt{║◌◌H◌◌◌◌║ 1} \\
\texttt{║◌◌◌◌◌◌◌║ 2} \\
\texttt{║◌◌◌◌◌┌─╢ 3} \\
\texttt{║◌IC┌─┤◌║ 4} \\
\texttt{║◌◌◌│W├─╢ 5} \\
\texttt{╟───┼─┘◌║ 6} \\
\texttt{║◌◌T│◌◌◌║ 7} \\
\texttt{╚═══╧═══╝} \\
\texttt{```} \\
\texttt{You are currently playing round 19 of maximum 20.} \\
\texttt{You have collectively accumulated 9 of 12 penalties.} \\
\texttt{L'altro giocatore ha spostato un oggetto sulla sua griglia.} \\
\\ 
\\ 
\texttt{Qual è il tuo prossimo comando?} \\
            }
        }
    }
    & \\ \\

    \theutterance \stepcounter{utterance}  
    & & & \multicolumn{4}{p{0.6\linewidth}}{
        \cellcolor[rgb]{0.9,0.9,0.9}{
            \makecell[{{p{\linewidth}}}]{
                \texttt{\tiny{[GM$\langle$P2]}}
                \texttt{COMUNICA: Sembra che gli oggetti 'C' siano stati allineati correttamente nella posizione (3,4). Ora che gli oggetti 'C' sono nella posizione corretta, penso che abbiamo fatto grandi progressi nel completare il gioco. Vorrei verificare se tutti gli oggetti sono nella posizione corretta e se è possibile completare il gioco con successo. Come stanno gli oggetti sulla tua griglia? Sono tutti nella posizione corretta?} \\
            }
        }
    }
     \\ \\

    \theutterance \stepcounter{utterance}  
    & & & \multicolumn{2}{p{0.3\linewidth}}{
        \cellcolor[rgb]{0.9,0.9,0.9}{
            \makecell[{{p{\linewidth}}}]{
                \texttt{\tiny{[GM$|$GM]}}
                \texttt{COMUNICA: Sembra che gli oggetti 'C' siano stati allineati correttamente nella posizione (3,4). Ora che gli oggetti 'C' sono nella posizione corretta, penso che abbiamo fatto grandi progressi nel completare il gioco. Vorrei verificare se tutti gli oggetti sono nella posizione corretta e se è possibile completare il gioco con successo. Come stanno gli oggetti sulla tua griglia? Sono tutti nella posizione corretta?} \\
            }
        }
    }
    & & \\ \\

    \theutterance \stepcounter{utterance}  
    & & \multicolumn{4}{p{0.6\linewidth}}{
        \cellcolor[rgb]{0.9,0.9,0.9}{
            \makecell[{{p{\linewidth}}}]{
                \texttt{\tiny{[P1$\langle$GM]}}
                \texttt{Spostato 'C' in (3,4) con successo. La tua griglia aggiornata appare così:} \\
\texttt{```} \\
\texttt{ 1234567} \\
\texttt{╔═══════╗} \\
\texttt{║◌◌H◌◌◌◌║ 1} \\
\texttt{║◌◌◌◌◌◌◌║ 2} \\
\texttt{║◌◌◌W◌┌─╢ 3} \\
\texttt{║◌◌C┌─┤◌║ 4} \\
\texttt{║◌I◌│◌├─╢ 5} \\
\texttt{╟───┼─┘◌║ 6} \\
\texttt{║◌◌T│◌◌◌║ 7} \\
\texttt{╚═══╧═══╝} \\
\texttt{```} \\
\texttt{You are currently playing round 19 of maximum 20.} \\
\texttt{You have collectively accumulated 9 of 12 penalties.} \\
\texttt{L'altro giocatore ha inviato questo messaggio:} \\
\\ 
\\ 
\texttt{"Sembra che gli oggetti 'C' siano stati allineati correttamente nella posizione (3,4). Ora che gli oggetti 'C' sono nella posizione corretta, penso che abbiamo fatto grandi progressi nel completare il gioco. Vorrei verificare se tutti gli oggetti sono nella posizione corretta e se è possibile completare il gioco con successo. Come stanno gli oggetti sulla tua griglia? Sono tutti nella posizione corretta?"} \\
\\ 
\\ 
\texttt{Qual è il tuo prossimo comando?} \\
            }
        }
    }
    & \\ \\

    \theutterance \stepcounter{utterance}  
    & \multicolumn{4}{p{0.6\linewidth}}{
        \cellcolor[rgb]{0.9,0.9,0.9}{
            \makecell[{{p{\linewidth}}}]{
                \texttt{\tiny{[P1$\rangle$GM]}}
                \texttt{COMUNICA: Sono felice di vedere che gli oggetti 'C' sono stati allineati correttamente! Dopo aver verificato la mia griglia, posso confermare che gli oggetti 'H' sono nella posizione (3,1), gli oggetti 'I' sono nella posizione (2,5), gli oggetti 'T' sono nella posizione (3,7) e gli oggetti 'W' sono nella posizione (4,3), ma dovrebbe essere nella posizione (5,5) per allinearsi con il tuo 'W'. Gli oggetti 'C' sono stati spostati nella posizione (3,4) come concordato. Sembra che gli oggetti 'W' siano gli unici che non sono ancora nella posizione corretta. Dovremmo spostare il mio 'W' nella posizione (5,5) per completare il gioco. Come pensi di procedere per completare il gioco?} \\
            }
        }
    }
    & & \\ \\

    \theutterance \stepcounter{utterance}  
    & & & \multicolumn{2}{p{0.3\linewidth}}{
        \cellcolor[rgb]{0.9,0.9,0.9}{
            \makecell[{{p{\linewidth}}}]{
                \texttt{\tiny{[GM$|$GM]}}
                \texttt{COMUNICA: Sono felice di vedere che gli oggetti 'C' sono stati allineati correttamente! Dopo aver verificato la mia griglia, posso confermare che gli oggetti 'H' sono nella posizione (3,1), gli oggetti 'I' sono nella posizione (2,5), gli oggetti 'T' sono nella posizione (3,7) e gli oggetti 'W' sono nella posizione (4,3), ma dovrebbe essere nella posizione (5,5) per allinearsi con il tuo 'W'. Gli oggetti 'C' sono stati spostati nella posizione (3,4) come concordato. Sembra che gli oggetti 'W' siano gli unici che non sono ancora nella posizione corretta. Dovremmo spostare il mio 'W' nella posizione (5,5) per completare il gioco. Come pensi di procedere per completare il gioco?} \\
            }
        }
    }
    & & \\ \\

    \theutterance \stepcounter{utterance}  
    & & \multicolumn{4}{p{0.6\linewidth}}{
        \cellcolor[rgb]{0.9,0.9,0.9}{
            \makecell[{{p{\linewidth}}}]{
                \texttt{\tiny{[GM$\rangle$P2]}}
                \texttt{Il tuo messaggio è stato inoltrato all'altro giocatore.} \\
\\ 
\\ 
\texttt{You are currently playing round 20 of maximum 20.} \\
\texttt{You have collectively accumulated 9 of 12 penalties.} \\
\texttt{L'altro giocatore ha inviato questo messaggio:} \\
\\ 
\\ 
\texttt{"Sono felice di vedere che gli oggetti 'C' sono stati allineati correttamente! Dopo aver verificato la mia griglia, posso confermare che gli oggetti 'H' sono nella posizione (3,1), gli oggetti 'I' sono nella posizione (2,5), gli oggetti 'T' sono nella posizione (3,7) e gli oggetti 'W' sono nella posizione (4,3), ma dovrebbe essere nella posizione (5,5) per allinearsi con il tuo 'W'. Gli oggetti 'C' sono stati spostati nella posizione (3,4) come concordato. Sembra che gli oggetti 'W' siano gli unici che non sono ancora nella posizione corretta. Dovremmo spostare il mio 'W' nella posizione (5,5) per completare il gioco. Come pensi di procedere per completare il gioco?"} \\
\\ 
\\ 
\texttt{Qual è il tuo prossimo comando?} \\
            }
        }
    }
    & \\ \\

    \theutterance \stepcounter{utterance}  
    & & & \multicolumn{4}{p{0.6\linewidth}}{
        \cellcolor[rgb]{0.9,0.9,0.9}{
            \makecell[{{p{\linewidth}}}]{
                \texttt{\tiny{[GM$\langle$P2]}}
                \texttt{COMUNICA: finito?} \\
            }
        }
    }
     \\ \\

    \theutterance \stepcounter{utterance}  
    & & & \multicolumn{2}{p{0.3\linewidth}}{
        \cellcolor[rgb]{0.9,0.9,0.9}{
            \makecell[{{p{\linewidth}}}]{
                \texttt{\tiny{[GM$|$GM]}}
                \texttt{COMUNICA: finito?} \\
            }
        }
    }
    & & \\ \\

    \theutterance \stepcounter{utterance}  
    & & & \multicolumn{2}{p{0.3\linewidth}}{
        \cellcolor[rgb]{0.9,0.9,0.9}{
            \makecell[{{p{\linewidth}}}]{
                \texttt{\tiny{[GM$|$GM]}}
                \texttt{Maximum number of rounds reached} \\
            }
        }
    }
    & & \\ \\

    \theutterance \stepcounter{utterance}  
    & & & \multicolumn{2}{p{0.3\linewidth}}{
        \cellcolor[rgb]{0.9,0.9,0.9}{
            \makecell[{{p{\linewidth}}}]{
                \texttt{\tiny{[GM$|$GM]}}
                \texttt{* success: True} \\
\texttt{* lose: False} \\
\texttt{* aborted: False} \\
\texttt{{-}{-}{-}{-}{-}{-}{-}} \\
\texttt{* Shifts: 4.00} \\
\texttt{* Max Shifts: 8.00} \\
\texttt{* Min Shifts: 4.00} \\
\texttt{* End Distance Sum: 3.24} \\
\texttt{* Init Distance Sum: 18.69} \\
\texttt{* Expected Distance Sum: 20.95} \\
\texttt{* Penalties: 9.00} \\
\texttt{* Max Penalties: 12.00} \\
\texttt{* Rounds: 20.00} \\
\texttt{* Max Rounds: 20.00} \\
\texttt{* Object Count: 5.00} \\
            }
        }
    }
    & & \\ \\

    \theutterance \stepcounter{utterance}  
    & & & \multicolumn{2}{p{0.3\linewidth}}{
        \cellcolor[rgb]{0.9,0.9,0.9}{
            \makecell[{{p{\linewidth}}}]{
                \texttt{\tiny{[GM$|$GM]}}
                \texttt{Initial Grids:} \\
\texttt{```} \\
\texttt{Player 1:     Player 2:} \\
\texttt{ 1234567       1234567} \\
\texttt{╔═══════╗     ╔═══════╗} \\
\texttt{║I◌◌◌◌◌◌║ 1   ║◌◌H◌◌◌◌║ 1} \\
\texttt{║◌◌◌◌◌◌◌║ 2   ║◌◌◌C◌◌◌║ 2} \\
\texttt{║◌◌◌T◌┌─╢ 3   ║◌◌◌◌◌┌─╢ 3} \\
\texttt{║◌◌◌┌─┤◌║ 4   ║◌◌◌┌─┤◌║ 4} \\
\texttt{║◌◌W│C├─╢ 5   ║◌IW│◌├─╢ 5} \\
\texttt{╟───┼─┘H║ 6   ╟───┼─┘◌║ 6} \\
\texttt{║◌◌◌│◌◌◌║ 7   ║T◌◌│◌◌◌║ 7} \\
\texttt{╚═══╧═══╝     ╚═══╧═══╝} \\
\texttt{```} \\
\\ 
\texttt{Final Grids:} \\
\texttt{```} \\
\texttt{Player 1:     Player 2:} \\
\texttt{ 1234567       1234567} \\
\texttt{╔═══════╗     ╔═══════╗} \\
\texttt{║◌◌H◌◌◌◌║ 1   ║◌◌H◌◌◌◌║ 1} \\
\texttt{║◌◌◌◌◌◌◌║ 2   ║◌◌◌◌◌◌◌║ 2} \\
\texttt{║◌◌◌W◌┌─╢ 3   ║◌◌◌◌◌┌─╢ 3} \\
\texttt{║◌◌C┌─┤◌║ 4   ║◌IC┌─┤◌║ 4} \\
\texttt{║◌I◌│◌├─╢ 5   ║◌◌◌│W├─╢ 5} \\
\texttt{╟───┼─┘◌║ 6   ╟───┼─┘◌║ 6} \\
\texttt{║◌◌T│◌◌◌║ 7   ║◌◌T│◌◌◌║ 7} \\
\texttt{╚═══╧═══╝     ╚═══╧═══╝} \\
\texttt{```} \\
            }
        }
    }
    & & \\ \\

\end{supertabular}
}

\end{document}
